%% 第22回編集委員会（物理第2巻）：組 熱W4
%% p2-06.tex  熱-6　もとにもどらない変化：エントロピー（★3〜5・3段）
%% 中身は tools/phys2-plan.md の B「熱-6」。公式の持ち主は A（N38 の ΔS＞0・N40・N48）。記号は C-1。
%% 数値はすべて Python で計算した：checks/物理第2巻/熱-6/calc_p2_06.py（出力 calc_p2_06_out.txt）
%% 水の比熱は本文 4.2 kJ/(kg·K)。3けたで書く例題（例題5）だけ 4.18 を使う。氷の融解熱 334 kJ/kg（熱-1 と同じ値〔巻末の表で照合〕）
%% 数学 8-7⑩ のコラム「閉じた道の面積は（物理 熱-6）」は ⑧ で受ける（手引き2章の10）。数学の本の文・数は変えない。
%% 数学にない道具「積分因子」は ⑥ で作る（計画書 C-6-2）。
%% 出典（本文で使った事実）
%%   ・クラウジウスが「エントロピー」の名前をつけた（1865年、ギリシャ語の「変化」tropē から。エネルギーに似せた）：
%%       R. Clausius, Annalen der Physik 125, 353 (1865)。Wikipedia "Entropy"・"Rudolf Clausius" で確かめた
%%   ・カルノー『火の動力についての考察』（1824年）の水車のたとえ：Wikipedia "Reflections on the Motive Power of Fire"
%%   ・失われる仕事＝T0×生まれるエントロピー：グイ（1889年）とストドラ（1905年ごろ）。Wikipedia "Gouy–Stodola theorem"
%%   ・「エクセルギー」の名前：ラント（Z. Rant）、1956年。Wikipedia "Exergy"
%%   ・時間の向きと宇宙のはじめの低いエントロピー（過去の仮説）：R. Penrose, The Emperor's New Mind (1989)、
%%       Wikipedia "Past hypothesis"。「なぜそうだったかは説明されていない」と書くにとどめる
%%   ・オンサーガーの相反定理（1931年、1968年ノーベル化学賞）
\makeatletter
% コップ（横から）  \qxgcup{(下のまん中)}{はば}{高さ}{水の色}
\newcommand\qxgcup[4]{\begin{scope}[shift={#1}]
  \fill[#4] ({-#2/2+0.55},0.5) -- ({#2/2-0.55},0.5) -- ({#2/2-0.55+0.25*#3*0.15},{#3*0.72}) -- ({-#2/2+0.55-0.25*#3*0.15},{#3*0.72}) -- cycle;
  \draw[bkInk, line width=0.7pt] ({-#2/2-0.1*#3*0.4},#3) -- ({-#2/2+0.4},0) -- ({#2/2-0.4},0) -- ({#2/2+0.1*#3*0.4},#3);\end{scope}}
% 区切りのある箱  \qxgbox{(左下)}{はば}{高さ}
\newcommand\qxgbox[3]{\draw[bkInk, line width=0.8pt] #1 rectangle ++(#2,#3);}
% 分子の点
\newcommand\qxgdot[1]{\fill[bkBlue] #1 circle (0.75);}
\makeatother

\phchapter{6}{もとにもどらない変化：\\エントロピー}{お湯と水をまぜたあと、なぜ元の温度に分かれない？}{\bklevelrange{3}{5}}{もとにもどらない変化：エントロピー}
\bkmichishirube{5-11,5-14,8-3,8-7,熱-2,熱-5}{熱-7,熱-8,熱-10,熱-15}

\begin{goal}
\begin{itemize}[label={\color{bkTheme}●}, leftmargin=1.6\zw, itemsep=0pt, topsep=1pt]
\item \bklevel{3}\enspace 熱÷温度 $\Delta S=Q/T$ を数えて、もどらない向きを言える。
\item \bklevel{4}\enspace 水・氷・気体の $\Delta S$ を計算し、$T$-$s$ 線図の面積を読める。
\item \bklevel{5}\enspace 生まれるエントロピーとエクセルギーで、機械の損失を数えられる。
\end{itemize}
\end{goal}

\phkoushiki{%
\phk{エントロピーの変化 $\Delta S=\dfrac{Q}{T}$}{高校物理（発展）}{ph:熱-6:dsqt}
\phk{自由膨張・混合でエントロピーが増える}{高校物理（発展）}{ph:熱-6:jiyuu}
}

\begin{hatena}
\noindent\begin{minipage}[c]{0.56\linewidth}\raggedright
80\,℃のお湯\qty{1}{kg}と20\,℃の水\qty{1}{kg}を混ぜると、50\,℃のぬるま湯になります。このぬるま湯が、ひとりでに80\,℃と20\,℃に分かれることはあるでしょうか。分かれても、エネルギーは保存しているのに。
\end{minipage}\hfill
\begin{minipage}[c]{0.42\linewidth}\centering
\begin{tikzpicture}[x=0.86mm, y=0.86mm]
  \qxgcup{(6,16)}{11}{13}{bkRed!35}
  \qxgcup{(22,16)}{11}{13}{bkBlue!30}
  \node[font=\sffamily\footnotesize] at (6,13) {80\,℃};
  \node[font=\sffamily\footnotesize] at (22,13) {20\,℃};
  \draw[bk arrow, line width=1pt] (31,23) -- (40,23);
  \qxgcup{(52,14)}{17}{17}{bkPurple!25}
  \node[font=\sffamily\footnotesize] at (52,11) {50\,℃};
  \draw[bk arrow, line width=1pt, bkGray] (40,5) -- (31,5);
  \node[font=\sffamily\bfseries\large, text=bkRed] at (35.5,9.5) {?};
  \node[font=\sffamily\footnotesize, text=bkGray] at (14,5) {もどる？};
\end{tikzpicture}
\end{minipage}
\end{hatena}

\noindent 熱-2では、熱と仕事を、同じエネルギーとして数えました（第一法則）。熱-5では、熱を全部仕事に変えることはできず、熱は自然には低い温度から高い温度へ移らないこと（第二法則）を学びました。この章では、この「自然に進む向き」を、1つの数で表します。その数が\textbf{エントロピー}です。エネルギーが「どれだけあるか」を数える量なら、エントロピーは「どちらへ進めるか」を数える量です。

\begin{yattemiyou}[インクのしずくは、もとにもどる？]
\textbf{用意するもの}\par
透明なコップ、水、インク（または食紅をとかした水）、スポイト、スマホのカメラ\par\vspace{2pt}
\noindent\begin{minipage}[c]{0.62\linewidth}
\begin{enumerate}[label={\bkcircled[bkGreen]{\arabic*}}, leftmargin=2.4\zw, itemsep=1pt, topsep=0pt]
\item コップに水を入れ、しばらく置いて、水の動きが止まるのを待つ。
\item スポイトで、インクを1滴、静かに落とす。かき混ぜない。
\item 10秒ごとに写真をとる（5分くらい）。動画でもよい。
\item 動画を逆に再生して、「ありえない」と感じる所をさがそう。
\end{enumerate}
\end{minipage}\hfill
\begin{minipage}[c]{0.35\linewidth}\centering
\begin{tikzpicture}[x=1mm, y=1mm]
  \qxgcup{(8,0)}{13}{20}{bkBlue!10}
  \fill[bkPurple!70] (8,12) circle (1.3);
  \draw[bk arrow] (16.5,10) -- (22.5,10);
  \qxgcup{(31,0)}{13}{20}{bkBlue!10}
  \fill[bkPurple!25] (31,8) ellipse (4.6 and 5.4);
  \fill[bkPurple!12] (31,7) ellipse (5.3 and 6.5);
\end{tikzpicture}
\end{minipage}
\end{yattemiyou}

\noindent やってみようの答え：インクは、ゆっくり広がって、やがてコップ全体が薄い色になります。逆に再生した動画では、薄く広がっていた色が、ひとりでに1か所に集まって、1滴にもどっていきます。こんなことは、何時間待っても起きません。でも、インクの粒1つ1つの動きは、ニュートンの運動方程式（力-2）にしたがっていて、その式は、時間を逆にしても同じ形のまま成り立ちます。1つ1つは逆に動けるのに、全体はもどらない。このふしぎを数で表すのが、この章です。

\begin{omoidasou}
\bkrecall{5-14}{$\ln x$ は $e$ を底とする対数。$\displaystyle\int\frac{1}{x}\,\dd x=\ln x+C$（$x>0$）。$\ln\dfrac{a}{b}=\ln a-\ln b$。}
\bkrecall{5-11}{定積分は「細い長方形の和」を細かくした先。$\displaystyle\int_a^bf(x)\,\dd x=F(b)-F(a)$。}
\bkrecall{8-7}{$p$-$V$ 図の上の道の線積分 $\int p\,\dd V$ は、気体が外にする仕事で、道によって変わる（⑩）。}
\bkrecall{熱-2}{第一法則 $\Delta U=Q+W$（$Q$ は受けとる熱、$W$ はされる仕事）。小さな変化では $\dd U=\delta Q+\delta W$。状態量の変化は全微分、熱と仕事は全微分ではない（⑦）。}
\bkrecall{熱-3}{理想気体の等温変化で、受けとる熱は $Q=nRT\ln\dfrac{V_2}{V_1}$（⑤）。$pV=nRT$。}
\bkrecall{熱-5}{熱は自然には低温から高温へ移らない。1つの温度の熱だけから仕事は取り出せない（③）。可逆な機関では $\dfrac{Q_{\mathrm L}}{Q_{\mathrm H}}=\dfrac{T_{\mathrm L}}{T_{\mathrm H}}$（⑤）。}
\end{omoidasou}

%% ======================== 物理の道具 ========================
\bksection{doguu}{物理の道具}

\phstage{me}{もどらない変化と、熱÷温度}{3}

\bksubsection{1}{もどらない変化の例}
\label{ph:熱-6:sec1}
身のまわりには、ひとりでは逆向きに起きない変化がたくさんあります。4つの例を見ましょう。

\begin{bkfigure}
\begin{tikzpicture}[x=1mm, y=1mm]
  % (ア) 混ぜる
  \qxgcup{(5,8)}{7}{9}{bkRed!35}\qxgcup{(14,8)}{7}{9}{bkBlue!30}
  \draw[bk arrow] (9.5,4) -- (9.5,0.5);
  \qxgcup{(9.5,-12)}{11}{10}{bkPurple!25}
  \node[font=\sffamily\bfseries\footnotesize, anchor=south] at (9.5,19.5) {（ア）混ぜる};
  % (イ) こする
  \begin{scope}[shift={(33,0)}]
  \fill[bkGray!35] (-3,-12) rectangle (25,-10);
  \filldraw[fill=bkOrange!35, draw=bkOrange!70!black, line width=0.6pt] (3,-10) rectangle (13,-4);
  \draw[bk arrow] (14,-7) -- (21,-7);
  \foreach \qxgx in {4,8,12} {\draw[bkRed, line width=0.6pt] (\qxgx,-3) .. controls (\qxgx+1,-1.5) and (\qxgx-1,-0.5) .. (\qxgx,1);}
  \node[font=\sffamily\footnotesize, align=center] at (11,9) {滑って止まり、\\熱になる};
  \node[font=\sffamily\bfseries\footnotesize, anchor=south] at (11,19.5) {（イ）擦る};
  \end{scope}
  % (ウ) 熱が流れる
  \begin{scope}[shift={(66,0)}]
  \filldraw[fill=bkRed!40, draw=bkRed!70!black, line width=0.6pt] (0,-8) rectangle (10,2);
  \filldraw[fill=bkBlue!25, draw=bkBlue!70!black, line width=0.6pt] (10,-8) rectangle (20,2);
  \draw[bk arrow, line width=1pt] (4,-3) -- (16,-3);
  \node[font=\sffamily\footnotesize] at (5,-11) {熱い};\node[font=\sffamily\footnotesize] at (15,-11) {冷たい};
  \node[font=\sffamily\footnotesize, align=center] at (10,9) {熱は高温から\\低温へ};
  \node[font=\sffamily\bfseries\footnotesize, anchor=south] at (10,19.5) {（ウ）熱が流れる};
  \end{scope}
  % (エ) ふくらむ
  \begin{scope}[shift={(98,0)}]
  \qxgbox{(0,-10)}{26}{12}
  \draw[bkInk, line width=0.8pt, dashed] (13,-10) -- (13,2);
  \foreach \qxgx/\qxgy in {2/-8,4/-3,6/-6,8/-1,10/-8,11/-4,3/0,7/-9,22/-7,17/-2} {\qxgdot{(\qxgx,\qxgy)}}
  \node[font=\sffamily\footnotesize, align=center] at (13,9) {仕切りをとると\\広がる};
  \node[font=\sffamily\bfseries\footnotesize, anchor=south] at (13,19.5) {（エ）ふくらむ};
  \end{scope}
\end{tikzpicture}
\bkfigcaption{もどらない変化の4つの例。どれも、矢印と逆の向きには、ひとりでは起きない。}\label{ph:熱-6:fig-rei}
\end{bkfigure}

\begin{enumerate}[label={\bkcircled{\arabic*}}, leftmargin=2.4\zw, itemsep=1pt, topsep=2pt]
\item （ア）お湯と水を混ぜると、ぬるま湯になる。ぬるま湯は、ひとりで熱い所と冷たい所に分かれない（\bkhatenaref）。
\item （イ）床の上で滑らせた箱は、摩擦で止まり、運動エネルギーは熱になる（力-4 ⑦、熱-2）。床が冷えて、その熱で箱が動き出すことはない。
\item （ウ）熱い物と冷たい物をくっつけると、熱は熱いほうから冷たいほうへ流れ、温度がそろう（熱-1 ①）。逆は起きない（熱-5 ③）。
\item （エ）箱の片側の気体は、仕切りをとると箱全体に広がる（自由膨張、熱-2 ⑥）。全部の分子が、ひとりで片側にもどることはない。
\end{enumerate}

4つとも、\textbf{エネルギーは保存しています}。（ア）では、お湯が失った熱と、水が受けとった熱は同じです。（イ）では、運動エネルギーが同じ大きさの熱になっただけです。だから、「逆向きに起きない」理由は、第一法則（エネルギーの保存）では説明できません。もう1つ、別の物差しがいります。

\phboxhead{可逆と不可逆}
ある変化のあと、\textbf{まわりに何の変化も残さずに}、もとの状態にもどせるとき、その変化を\textbf{可逆}な変化といいます。\bkidx{かぎゃく}{可逆変化}もどせないとき、\textbf{不可逆}な変化といいます。4つの例は、どれも不可逆です。（ア）のぬるま湯も、冷蔵庫と電気ストーブを使えば、半分を80\,℃、半分を20\,℃にもどせます。でも、そのとき電気を使い、部屋に熱を捨てるので、「まわりに変化が残る」のです。

\begin{seikaku}
本当の変化は、どれも少しずつ不可逆です。可逆な変化は、摩擦がなく、温度や圧力の差がほとんど0のまま、とてもゆっくり進む変化（準静的な変化、熱-2 ③）の、行きつく先として考える「理想」です。力学で摩擦のない面を考えるのと同じです。
\end{seikaku}

%% ---------------------------------------------------------------
\bksubsection{2}{熱÷温度を数える：$\Delta S=Q/T$}
\label{ph:熱-6:sec2}
\phboxhead{可逆な機関を通りぬけても変わらない量}
熱-5 ⑤のカルノーの機関（可逆な熱機関）を思い出しましょう。高温 $T_{\mathrm H}$ の所から熱 $Q_{\mathrm H}$ を受けとり、低温 $T_{\mathrm L}$ の所へ熱 $Q_{\mathrm L}$ を捨て、差を仕事 $W_{\text{出}}$ として取り出します（どれも大きさで、正の数。熱-5 と同じ書き方）。可逆な機関では、
\[\frac{Q_{\mathrm L}}{Q_{\mathrm H}}=\frac{T_{\mathrm L}}{T_{\mathrm H}}\qquad\text{だから}\qquad \frac{Q_{\mathrm H}}{T_{\mathrm H}}=\frac{Q_{\mathrm L}}{T_{\mathrm L}}\]
でした。熱 $Q$ は、機関を通ると仕事の分だけ減ります（$Q_{\mathrm L}<Q_{\mathrm H}$）。ところが、\textbf{熱を温度でわった量 $Q/T$ は、入るときと出るときで同じ}です。

\begin{bkfigure}
\begin{tikzpicture}[x=1mm, y=1mm]
  \node[bk node, minimum width=30mm, minimum height=8mm, fill=bkRed!15, font=\sffamily\small] (h) at (0,38) {高温 $T_{\mathrm H}=\qty{600}{K}$};
  \node[bk node, minimum width=18mm, minimum height=10mm, font=\sffamily\small, align=center] (e) at (0,15) {可逆な\\機関};
  \node[bk node, minimum width=30mm, minimum height=8mm, fill=bkBlue!12, font=\sffamily\small] (l) at (0,-8) {低温 $T_{\mathrm L}=\qty{300}{K}$};
  \draw[bk link, line width=1.6pt] (h) -- node[right=1.5mm, font=\sffamily\footnotesize, align=left] {$Q_{\mathrm H}=\qty{3460}{J}$\\$\div600=\qty{5.76}{J/K}$} (e);
  \draw[bk link, line width=1.1pt] (e) -- node[right=1.5mm, font=\sffamily\footnotesize, align=left] {$Q_{\mathrm L}=\qty{1730}{J}$\\$\div300=\qty{5.76}{J/K}$} (l);
  \draw[bk link, bkOrange, line width=1.1pt] (e.west) -- ++(-10,0) node[left, font=\sffamily\footnotesize, text=bkInk] {$W_{\text{出}}=\qty{1730}{J}$};
  % 右：水車
  \begin{scope}[shift={(66,0)}]
  \fill[bkBlue!20] (-22,26) rectangle (-4,30);
  \draw[bkBlue!70!black, line width=0.6pt] (-22,26) -- (-4,26);
  \draw[bkBlue!60, line width=2pt] (-5,27.5) .. controls (0,27.5) and (1,24) .. (1,20);
  \draw[bkInk, line width=0.8pt] (6,12) circle (8);
  \foreach \qxga in {0,45,...,315} {\draw[bkInk, line width=0.7pt] (6,12) -- ++(\qxga:8);}
  \draw[bkBlue!60, line width=2pt] (11,5.5) .. controls (13,2) and (14,0) .. (16,-2);
  \fill[bkBlue!20] (8,-6) rectangle (30,-2);
  \draw[bkBlue!70!black, line width=0.6pt] (8,-2) -- (30,-2);
  \node[font=\sffamily\footnotesize, align=center, anchor=north] at (6,-7.5) {カルノーの水車のたとえ\\（出る水の量は、入る量と同じ）};
  \end{scope}
\end{tikzpicture}
\bkfigcaption{熱-5 例題8 のカルノーの機関（理想気体 \qty{1}{mol}）。熱 $Q$ は減るが、$Q/T$ は同じ \qty{5.76}{J/K} のまま通りぬける。水車を回した水が、高さを失っても量は減らないのと似ている。}\label{ph:熱-6:fig-carnot}
\end{bkfigure}

カルノーは1824年の本で、熱機関を、高い所から低い所へ落ちる水で回る水車にたとえました。水は高さ（エネルギー）を失いますが、量は減りません。今のことばで言えば、可逆な機関を通りぬけても量が変わらないのは、熱 $Q$ ではなく、$Q/T$ のほうです。

\phboxhead{エントロピーの変化を決める}
\label{ph:熱-6:dsqt}そこで、温度 $T$ の物が、温度を変えずに、可逆に熱 $Q$ を受けとったとき、その物の\textbf{エントロピー} $S$ が
\pheq{\Delta S=\frac{Q}{T}}{J/K ＝ J ÷ K}
だけ増える、と決めます。\bkidx{えんとろぴー}{エントロピー}熱を出したときは $Q$ が負なので、$\Delta S$ も負（減る）です。単位は J/K です。$T$ は必ず絶対温度（K）で入れます。

\begin{phyomi}[$\Delta S=\dfrac{Q}{T}$]
「受けとった熱を、そのときの絶対温度でわる」。同じ熱でも、冷たい物が受けとるほど、エントロピーは大きく増える。熱い物にとっての\qty{100}{J}より、冷たい物にとっての\qty{100}{J}のほうが「大きなできごと」だ、と読むこともできる。
\end{phyomi}

\begin{phmijika}[お風呂と冷蔵庫]
\qty{300}{K}（約27\,℃）の大きな風呂おけの湯に、\qty{600}{J}の熱が入ると、湯の温度はほとんど変わらず、エントロピーは $600\div300=\qty{2.0}{J/K}$ 増えます。同じ\qty{600}{J}が、\qty{250}{K}（約 $-23$\,℃）の冷凍庫の中の大きな氷のかたまりに入れば、$600\div250=\qty{2.4}{J/K}$ 増えます。
\end{phmijika}

\begin{phbaisuru}[半分]{温度 $T$}
同じ熱 $Q$ を受けとっても、$T$ を半分にすると $\Delta S$ は2倍になります（反比例、3-7）。
\par\smallskip\noindent\begin{tikzpicture}
\begin{axis}[bkplot, width=96mm, height=42mm, xmin=0, xmax=650, ymin=0, ymax=6.4, xtick={150,300,600}, ytick={1,2,4}, xlabel={$T$（K）}, ylabel={$\Delta S$（J/K）}]
  \addplot[bkcurve, domain=100:640] {600/x};
  \fill[bkRed] (axis cs:150,4) circle (1.8pt) (axis cs:300,2) circle (1.8pt) (axis cs:600,1) circle (1.8pt);
  \node[font=\sffamily\footnotesize, anchor=south west, text=bkRed!75!black] at (axis cs:305,2.2) {$T$ を半分にすると2倍};
  \node[font=\sffamily\footnotesize, anchor=north east, text=bkGray] at (axis cs:640,6.2) {$Q=\qty{600}{J}$ のとき};
\end{axis}
\end{tikzpicture}
\end{phbaisuru}

\phboxhead{熱が流れると、エントロピーの合計は増える}
熱い物（$T_{\mathrm H}$）から冷たい物（$T_{\mathrm L}$）へ、熱 $Q$ が流れたとします。どちらの物も大きくて、温度はほとんど変わらないとします。熱い物は $Q$ を出したので $-Q/T_{\mathrm H}$、冷たい物は $Q$ を受けとったので $+Q/T_{\mathrm L}$ だけ、エントロピーが変わります。合計は
\[\Delta S_{\text{合計}}=-\frac{Q}{T_{\mathrm H}}+\frac{Q}{T_{\mathrm L}}=Q\Bigl(\frac{1}{T_{\mathrm L}}-\frac{1}{T_{\mathrm H}}\Bigr)>0\]
です（$T_{\mathrm L}<T_{\mathrm H}$ なので $1/T_{\mathrm L}>1/T_{\mathrm H}$）。熱が自然に流れる向きでは、エントロピーの合計は\textbf{必ず増えます}。もし逆に、冷たい物から熱い物へ熱がひとりで流れたら、合計は減ります。熱-5 ③で学んだとおり、それは起きません。

ここから、次のように言えそうです。
\begin{taisetsu}
\textbf{エントロピー増大の法則}：まわりと熱も仕事もやりとりしない（孤立した）物の集まりでは、エントロピーの合計は、増えるか、変わらないかのどちらかで、減ることはない。可逆な変化なら変わらず、不可逆な変化なら増える。
\end{taisetsu}
これは第二法則の、3つめの言い方です。熱が流れる場合で確かめましたが、どんな変化でも成り立つことを、⑥で熱-5の第二法則から導きます。\bkidx{えんとろぴーぞうだいのほうそく}{エントロピー増大の法則}

\begin{seikaku}
$\Delta S=Q/T$ の $Q$ は、\textbf{可逆に}受けとった熱です。熱が流れる例では、熱い物も冷たい物も、それぞれの中では温度がそろったまま、ゆっくり熱を出し入れしたと考えています（不可逆なのは、温度の違う2つの物の\textbf{境目}で、熱が流れる所だけ）。どんな場合に $Q/T$ が使えるかは、⑥で正確に決めます。
\end{seikaku}

\Needspace{16\baselineskip}\phstage{te}{表計算で、熱の移動を少しずつ追う}{3}

\bksubsection{3}{お湯から水へ、少しずつ熱を移す}
\label{ph:熱-6:sec3}
\bkhatenaref のお湯（80\,℃、\qty{1}{kg}）と水（20\,℃、\qty{1}{kg}）を、薄い金属の板ごしにくっつけて、熱がゆっくり流れるようにします。こんどは、どちらの温度も変わっていくので、$\Delta S=Q/T$ をそのまま1回では使えません。そこで、\textbf{少しずつに分けて}数えます。

\begin{enumerate}[label={\bkcircled{\arabic*}}, leftmargin=2.4\zw, itemsep=1pt, topsep=2pt]
\item 1回に、お湯から水へ $\Delta Q=4200\times1=\qty{4200}{J}$ の熱を移す。お湯は\qty{1}{K}下がり、水は\qty{1}{K}上がる（$Q=mc\Delta T$、熱-1 ③）。
\item 1回の間に、お湯の温度は少し変わるので、その回の\textbf{まん中の温度}でわる。1回めなら、お湯は79.5\,℃（\qty{352.65}{K}）、水は20.5\,℃（\qty{293.65}{K}）。
\item お湯の変化 $-\Delta Q/T_{\text{お湯}}$ と、水の変化 $+\Delta Q/T_{\text{水}}$ をたして、その回の合計の増えを出す。
\item 温度がそろう（50\,℃）まで、30回繰り返して、増えを足していく。
\end{enumerate}

\begin{bkcenter}
\bksheet{7}{13.5}{%
1 & 回 & お湯（℃） & 水（℃） & お湯の $\Delta S$ & 水の $\Delta S$ & 増え & 合計 \\
2 & 1 & 79 & 21 & $-11.910$ & 14.303 & 2.393 & 2.39 \\
3 & 2 & 78 & 22 & $-11.944$ & 14.254 & 2.311 & 4.70 \\
4 & 3 & 77 & 23 & $-11.978$ & 14.206 & 2.228 & 6.93 \\
5 & $\vdots$ & $\vdots$ & $\vdots$ & $\vdots$ & $\vdots$ & $\vdots$ & $\vdots$ \\
6 & 10 & 70 & 30 & $-12.222$ & 13.877 & 1.656 & 20.24 \\
7 & 20 & 60 & 40 & $-12.588$ & 13.434 & 0.846 & 32.33 \\
8 & 28 & 52 & 48 & $-12.897$ & 13.098 & 0.201 & 36.19 \\
9 & 29 & 51 & 49 & $-12.937$ & 13.058 & 0.121 & 36.31 \\
10 & 30 & 50 & 50 & $-12.977$ & 13.017 & 0.040 & 36.35 \\}
\bktabcaption{エントロピーの変化（J/K）を1回ずつ数える。D2＝$-$4200/(B2＋0.5＋273.15)、E2＝4200/(C2－0.5＋273.15)、F2＝D2＋E2、G3＝G2＋F3。}\label{ph:熱-6:tab-sheet}
\end{bkcenter}

\begin{bkfigure}
\begin{tikzpicture}
\begin{axis}[bkplot, width=112mm, height=50mm, xmin=0, xmax=32, ymin=0, ymax=42, xtick={10,20,30}, ytick={10,20,30,36.35}, yticklabels={10,20,30,36.4}, xlabel={回}, ylabel={エントロピーの合計の増え（J/K）}]
  \addplot[bkBlue, mark=*, only marks] coordinates {(1,2.39)(2,4.70)(3,6.93)(4,9.08)(5,11.14)(6,13.12)(7,15.02)(8,16.84)(9,18.58)(10,20.24)(11,21.81)(12,23.30)(13,24.72)(14,26.05)(15,27.30)(16,28.46)(17,29.55)(18,30.56)(19,31.49)(20,32.33)(21,33.10)(22,33.78)(23,34.38)(24,34.91)(25,35.35)(26,35.71)(27,35.99)(28,36.19)(29,36.31)(30,36.35)};
  \addplot[bkRed, dashed, line width=0.8pt, domain=0:32] {36.35};
  \node[font=\sffamily\footnotesize, anchor=south east, text=bkRed!75!black] at (axis cs:31,37) {温度がそろうと、増えが止まる};
  \node[font=\sffamily\footnotesize, anchor=north west, text=bkBlue!80!black] at (axis cs:12,18) {はじめは速く増える};
\end{axis}
\end{tikzpicture}
\bkfigcaption{表\ref{ph:熱-6:tab-sheet}の「合計」の列をグラフにした。いつも増え続け、温度がそろう30回めで止まる。}\label{ph:熱-6:fig-sheet}
\end{bkfigure}

\phboxhead{表から読みとれること}
\begin{enumerate}[label={\bkcircled{\arabic*}}, leftmargin=2.4\zw, itemsep=1pt, topsep=2pt]
\item お湯のエントロピーは毎回減り、水のエントロピーは毎回増えます。でも、水のほうが温度が低いので、同じ\qty{4200}{J}でも増え方が大きく、\textbf{合計は毎回増えます}。
\item 増えは、温度の差が大きいはじめほど大きく（2.39）、差が小さくなると0に近づきます（0.040）。温度がそろうと、熱が流れなくなり、増えも止まります。
\item 合計は\qty{36.35}{J/K}です。⑦で積分して求める値（\qty{36.355}{J/K}）と、ほとんど同じです。
\end{enumerate}

\begin{phmijika}[逆向きに1回だけ移すと]
もし、1回めに逆向きに、水（20\,℃）からお湯（80\,℃）へ\qty{4200}{J}を移すと、その回の増えは
\[\frac{4200}{352.65}-\frac{4200}{293.65}\approx\qty{-2.39}{J/K}\]
で、合計が減ります。②のとおり、これはひとりでは起きない向きです。
\end{phmijika}

\begin{seikaku}
1回で\qty{30}{K}分（\qty{126000}{J}）を一度に移して、はじめの温度（80\,℃と20\,℃）でわると、
\[-\frac{126000}{353.15}+\frac{126000}{293.15}\approx\qty{73}{J/K}\]
となり、2倍もずれます。温度が変わるときは、細かく分けて、そのときの温度でわる必要があります。分ける数をどんどん増やした先が、5-11の定積分です（⑦）。
\end{seikaku}

%% ---------------------------------------------------------------
\bksubsection{4}{自由膨張と混合のエントロピー}
\label{ph:熱-6:sec4}
\phboxhead{熱が出入りしないのに、エントロピーが増える？}
熱を通さない箱を仕切りで2つに分け、左半分（体積 $V$）に理想気体 $n$〔mol〕を入れ、右半分は真空にします。仕切りをとると、気体は箱全体（体積 $2V$）に広がります（自由膨張）。熱-2 ⑥で学んだとおり、このとき気体は熱を受けとらず（$Q=0$）、仕事もしない（真空へ広がるので、押す相手がいない）ので、$\Delta U=0$ で、理想気体の温度は変わりません。

では、$Q=0$ だから $\Delta S=0$ でしょうか。そうではありません。$\Delta S=Q/T$ の $Q$ は、\textbf{可逆に}受けとった熱でした。自由膨張は不可逆なので、この $Q=0$ を使ってはいけないのです。

\phboxhead{同じはじめとおわりを、可逆な道でつなぐ}
エントロピーは、\textbf{いまの状態（$p$・$V$・$T$）だけで決まる量}（状態量）です（⑥で確かめます）。だから、はじめの状態（体積 $V$、温度 $T$）とおわりの状態（体積 $2V$、温度 $T$）が同じなら、どんな道で変化させても、エントロピーの変化は同じです。そこで、自由膨張の代わりに、同じはじめとおわりを結ぶ\textbf{可逆な道}を考えます。

\begin{bkfigure}
\begin{tikzpicture}[x=1mm, y=1mm]
  % （ア）自由膨張
  \node[font=\sffamily\bfseries\footnotesize, anchor=south] at (55,19) {（ア）自由膨張（不可逆）：$Q=0$、$W=0$、$T$ は同じ};
  \qxgbox{(0,0)}{44}{16}
  \draw[bkInk, line width=1pt] (22,0) -- (22,16);
  \foreach \qxgx/\qxgy in {3/3,7/8,11/4,15/12,5/13,18/6,9/1.8,13/9} {\qxgdot{(\qxgx,\qxgy)}}
  \node[font=\sffamily\footnotesize] at (33,8) {真空};
  \draw[bk arrow, line width=1pt] (48,6) -- (62,6);
  \node[font=\sffamily\footnotesize, anchor=south, align=center] at (55,7) {仕切りを\\とる};
  \qxgbox{(66,0)}{44}{16}
  \foreach \qxgx/\qxgy in {70/3,78/12,84/5,93/10,101/3,74/8,98/13,106/7} {\qxgdot{(\qxgx,\qxgy)}}
  % （イ）可逆な道
  \begin{scope}[shift={(0,-36)}]
  \node[font=\sffamily\bfseries\footnotesize, anchor=south] at (55,19) {（イ）同じはじめとおわりを結ぶ可逆な道：等温でゆっくり膨張};
  \fill[bkRed!12] (-2,-5) rectangle (112,0);
  \node[font=\sffamily\footnotesize, anchor=north] at (55,-5.5) {温度 $T$ の大きな水そう（気体はここから熱 $Q$ を受けとる）};
  \qxgbox{(0,0)}{44}{16}
  \fill[bkGray!60] (21,0.5) rectangle (23,15.5);
  \foreach \qxgx/\qxgy in {3/3,7/8,11/4,15/12,5/13,18/6,9/1.8,13/9} {\qxgdot{(\qxgx,\qxgy)}}
  \draw[bk arrow, bkRed, line width=1pt] (10,-4) -- (10,-0.6);
  \draw[bk arrow, bkOrange] (24.5,8) -- (32,8);
  \draw[bk arrow, line width=1pt] (48,6) -- (62,6);
  \node[font=\sffamily\footnotesize, anchor=south, align=center] at (55,7) {ゆっくり\\押し返す};
  \qxgbox{(66,0)}{44}{16}
  \fill[bkGray!60] (107.5,0.5) rectangle (109.5,15.5);
  \foreach \qxgx/\qxgy in {70/3,78/12,84/5,93/10,101/3,74/8,98/13,104/7} {\qxgdot{(\qxgx,\qxgy)}}
  \draw[bk arrow, bkRed, line width=1pt] (88,-4) -- (88,-0.6);
  \end{scope}
\end{tikzpicture}
\bkfigcaption{（ア）の自由膨張と、（イ）の等温の可逆な膨張は、はじめとおわりの状態が同じ。だから気体のエントロピーの変化も同じ。}\label{ph:熱-6:fig-jiyuu}
\end{bkfigure}

\label{ph:熱-6:jiyuu}（イ）では、気体はピストンを押して外へ仕事をしながら、温度 $T$ の水そうから同じだけの熱を受けとり、温度を変えずに $V$ から $2V$ へ広がります。熱-3 ⑤の等温変化の式から、受けとる熱は
\[Q=nRT\ln\frac{2V}{V}=nRT\ln2\]
です。温度 $T$ は一定なので、②の $\Delta S=Q/T$ が使えて、
\pheq{\Delta S=\frac{nRT\ln2}{T}=nR\ln2}{J/K ＝ mol × J/(mol·K)}
です。\qty{1}{mol}なら $8.314\times0.693\approx\qty{5.76}{J/K}$ です。（ア）の自由膨張でも、気体のエントロピーは同じ\qty{5.76}{J/K}だけ増えます。

\phboxhead{可逆な道と自由膨張の違い}
（イ）では、水そうが熱 $Q$ を出したので、水そうのエントロピーは $-Q/T=\qty{-5.76}{J/K}$ です。気体と水そうの合計は0で、変わりません（可逆）。（ア）では、まわりは何も変わらないので、合計は $+\qty{5.76}{J/K}$ です（不可逆）。\textbf{気体だけを見れば同じ変化でも、まわりまで入れた合計で、可逆か不可逆かが分かれます}。

\begin{phbaisuru}[2乗]{体積の倍率}
体積を2倍にすると $nR\ln2$、4倍（2倍の2乗）にすると $nR\ln4=2nR\ln2$ です。体積の\textbf{倍率}を2乗にすると、エントロピーの増えは2倍になります（かけ算がたし算になる対数の性質、3-12③）。
\end{phbaisuru}

\phboxhead{2種類の気体が混ざる}
窒素\qty{1}{mol}と酸素\qty{1}{mol}を、同じ温度・同じ圧力で、仕切りの左右に入れます。仕切りをとると、2つの気体は混ざり合います。それぞれの気体は、箱全体に広がるので、どちらも体積が2倍になったのと同じです（理想気体の分子は、ほかの種類の分子を気にせず動くから）。合計は
\[\Delta S=R\ln2+R\ln2=2R\ln2\approx\qty{11.5}{J/K}\]
です。混ざった気体が、ひとりで窒素と酸素に分かれないのは、分かれるとエントロピーが減るからです。

\begin{phmijika}[分子の数で考えると]
分子が1個なら、右半分から左半分にもどってくることは、よくあります。2個とも左にいる確率は $\tfrac14$、10個なら約1000回に1回です。でも100個が全部左にいる確率は $\bigl(\tfrac12\bigr)^{100}\approx8\times10^{-31}$ です。\qty{1}{mol}は約 $6\times10^{23}$ 個なので、全部がもどることは、ありえないと言ってよいのです。エントロピーと「並び方の数」の関係は、熱-15 で学びます（$nR\ln2$ の $\ln2$ は、分子1個あたり「左か右か」の2通りから出てきます）。
\end{phmijika}

%% ---------------------------------------------------------------
\Needspace{22\baselineskip}\bksubsection{5}{はてな？の答え：お湯と水は、なぜ分かれない？}
\label{ph:熱-6:sec5}
\begin{phmodel}[お湯と水を混ぜる]
\phignore{なべ・まわりの空気との熱のやりとり。水の比熱の温度による変化。混ぜるときの水の運動エネルギー。}
\phassume{比熱は $c=\qty{4.2}{kJ/(kg.K)}$ で一定。混ぜたあと、全体の温度は1つにそろう。}
\phlimit{なべが熱をたくさん吸うとき。量がとても少なく、まわりへ逃げる熱が効くとき。}
\end{phmodel}

\bkhatenaref の答えは、③の表計算でわかりました。お湯と水が混ざって50\,℃になるとき、全体のエントロピーは約\qty{36}{J/K}増えます（⑦の式で計算すると\qty{36.4}{J/K}。例題1）。

\begin{enumerate}[label={\bkcircled{\arabic*}}, leftmargin=2.4\zw, itemsep=1pt, topsep=2pt]
\item ぬるま湯が、ひとりで80\,℃の半分と20\,℃の半分に分かれたとすると、それは③の逆なので、エントロピーの合計は\qty{36}{J/K}\textbf{減ります}。
\item まわりと何もやりとりしない（孤立した）ぬるま湯では、エントロピーの合計は減らない（②のエントロピー増大の法則）。だから、分かれることはありません。
\item エネルギーの保存（第一法則）だけなら、分かれてもかまいません。分かれない理由は、第二法則です。
\end{enumerate}

\phboxhead{まわりに変化を残せば、分けられる}
冷凍機（ヒートポンプ）を使えば、半分から熱をくみ出して冷やし、残りの半分を温めることはできます。そのときは、電気の仕事を使い、使った仕事の分の熱を部屋へ捨てます。部屋が受けとった熱で、部屋のエントロピーが増えるので、全体の合計は減りません。いちばんうまくやっても、部屋の温度を20\,℃として約\qty{11}{kJ}の仕事がいります（例題8。⑨）。\textbf{もどらない変化は、まわりに「エントロピーを捨てる」ことで、やっともどせる}のです。

\begin{phmada}[時間は、なぜ一方向にしか進まない？]
力-2の運動方程式も、電気や磁気の法則も、時間を逆にしても同じ形で成り立ちます。それなのに、私たちの世界では、インクは広がるだけで集まらず、エントロピーは増えるだけです。この「時間の向き」は、宇宙のはじめ（ビッグバンのころ）の状態が、とても低いエントロピーだったからだと考えられています。でも、なぜ宇宙がそんなに特別な状態から始まったのかは、まだわかっていません。ペンローズなどが問題にしてきた、宇宙論と物理の根元にかかわる問いで、今も議論が続いています。
\end{phmada}

\phkokomade

\Needspace{16\baselineskip}\phstage{shiki}{クラウジウスの不等式から導く}{5}

\bksubsection{6}{クラウジウスの不等式と、状態量 $S$}
\label{ph:熱-6:sec6}
②では、エントロピーを「温度一定で可逆に受けとった熱÷温度」と決め、熱が流れる場合だけで増大の法則を確かめました。ここでは、熱-5 ③の第二法則（1つの温度の熱だけから仕事は取り出せない）から出発して、どんな変化でも使える形にします。

\phboxhead{熱を「受けとる向きを正」にして数える}
この節では、熱 $Q$ を、働く物（エンジンの中の気体など）が\textbf{受けとる向きを正}にして書きます（熱-2 の決まり）。捨てる熱は負の数です。カルノーの機関なら、温度 $T_1=T_{\mathrm H}$ で $Q_1=Q_{\mathrm H}$、温度 $T_2=T_{\mathrm L}$ で $Q_2=-Q_{\mathrm L}$ です。②の可逆な機関の式は
\[\frac{Q_1}{T_1}+\frac{Q_2}{T_2}=\frac{Q_{\mathrm H}}{T_{\mathrm H}}-\frac{Q_{\mathrm L}}{T_{\mathrm L}}=0\]
です。不可逆な機関（摩擦や、温度の差のある所での熱の流れがある）では、効率がカルノーより低いので（熱-5 ⑤のカルノーの定理）、$Q_{\mathrm L}/Q_{\mathrm H}>T_{\mathrm L}/T_{\mathrm H}$、つまり
\[\frac{Q_1}{T_1}+\frac{Q_2}{T_2}<0\]
です（例題13）。

\phboxhead{たくさんの温度と熱をやりとりするサイクル}
こんどは、働く物が1周する間に、温度 $T_1$、$T_2$、…、$T_N$ のたくさんの熱の源（大きな水そうなど）から、熱 $\delta Q_1$、$\delta Q_2$、…、$\delta Q_N$ を受けとる場合を考えます。次の手順で、$\displaystyle\sum_i\frac{\delta Q_i}{T_i}\leqq0$ を示します。

\begin{bkfigure}
\begin{tikzpicture}[x=1mm, y=1mm]
  \node[bk node, minimum width=100mm, minimum height=7mm, fill=bkRed!8, font=\sffamily\small] (t0) at (40,34) {ただ1つの水そう $T_0$};
  \node[bk node, minimum width=15mm, minimum height=8mm, font=\sffamily\footnotesize] (c1) at (8,20) {可逆 $C_1$};
  \node[bk node, minimum width=15mm, minimum height=8mm, font=\sffamily\footnotesize] (c2) at (34,20) {可逆 $C_2$};
  \node[font=\sffamily\small] at (55,20) {…};
  \node[bk node, minimum width=15mm, minimum height=8mm, font=\sffamily\footnotesize] (cn) at (76,20) {可逆 $C_N$};
  \node[bk node, minimum width=15mm, minimum height=7mm, fill=bkOrange!12, font=\sffamily\footnotesize] (r1) at (8,6) {$T_1$};
  \node[bk node, minimum width=15mm, minimum height=7mm, fill=bkOrange!12, font=\sffamily\footnotesize] (r2) at (34,6) {$T_2$};
  \node[bk node, minimum width=15mm, minimum height=7mm, fill=bkOrange!12, font=\sffamily\footnotesize] (rn) at (76,6) {$T_N$};
  \node[bk node, minimum width=80mm, minimum height=8mm, fill=bkBlue!10, font=\sffamily\small] (g) at (42,-10) {1周する働く物（調べたいサイクル）};
  \foreach \qxga in {1,2,n} {\draw[bk link] (t0.south -| c\qxga) -- (c\qxga); \draw[bk link] (c\qxga) -- (r\qxga); \draw[bk link] (r\qxga) -- (g.north -| r\qxga);}
  \node[font=\sffamily\footnotesize, anchor=west] at (9.5,-1.8) {$\delta Q_1$};
  \node[font=\sffamily\footnotesize, anchor=west] at (35.5,-1.8) {$\delta Q_2$};
  \node[font=\sffamily\footnotesize, anchor=west] at (77.5,-1.8) {$\delta Q_N$};
  \node[font=\sffamily\footnotesize, anchor=west] at (9.5,27.5) {$T_0\,\delta Q_1/T_1$};
  \node[font=\sffamily\footnotesize, anchor=west] at (77.5,27.5) {$T_0\,\delta Q_N/T_N$};
  \draw[bk link, bkOrange] (g.east) -- ++(14,0) node[right, font=\sffamily\footnotesize, text=bkInk] {仕事};
\end{tikzpicture}
\bkfigcaption{クラウジウスの不等式の考え方。水そう $T_i$ が働く物に出した熱 $\delta Q_i$ を、可逆な機関 $C_i$ がちょうど補う。全体では、ただ1つの水そう $T_0$ とだけ熱をやりとりする。}\label{ph:熱-6:fig-clausius}
\end{bkfigure}

\begin{enumerate}[label={\bkcircled{\arabic*}}, leftmargin=2.4\zw, itemsep=1pt, topsep=2pt]
\item それぞれの水そう $T_i$ に、可逆なカルノーの機関 $C_i$ をつなぎ、別の大きな水そう $T_0$ から熱を受けとって、水そう $T_i$ へちょうど $\delta Q_i$ を渡すようにする。すると、水そう $T_i$ は、出した熱をすぐ補われるので、何も変わらない。
\item $C_i$ は可逆なので、②の式（$Q_{\mathrm H}/T_{\mathrm H}=Q_{\mathrm L}/T_{\mathrm L}$）から、$T_0$ から受けとる熱は $T_0\,\delta Q_i/T_i$。
\item 働く物と $C_1$〜$C_N$ をまとめて1つの機械と見ると、これは1周するたびにもとにもどり、熱は $T_0$ の水そうとだけやりとりする。受けとる熱の合計は $T_0\displaystyle\sum_i\frac{\delta Q_i}{T_i}$。
\item 1周で $\Delta U=0$ なので、第一法則から、この熱の合計は、まとめた機械が外へする仕事に等しい。
\item 1つの温度の熱だけから、正の仕事を取り出すことはできない（熱-5 ③）。だから、この仕事は0か負。$T_0>0$ なので、$\displaystyle\sum_i\frac{\delta Q_i}{T_i}\leqq0$。
\end{enumerate}
熱の源を細かくして、1周の積分で書くと、
\pheq{\oint\frac{\delta Q}{T}\leqq0}{J/K}
です。これを\textbf{クラウジウスの不等式}といいます。\bkidx{くらうじうすのふとうしき}{クラウジウスの不等式}$T$ は、熱 $\delta Q$ を渡す相手（熱の源）の温度です。

\phboxhead{可逆なサイクルでは0：状態量 $S$}
可逆なサイクルは、逆に回すこともできます。逆に回すと、すべての $\delta Q$ の符号が変わるので、$\oint\delta Q/T\geqq0$ にもなります。両方が成り立つのは、0のときだけです。
\[\oint\frac{\delta Q_{\mathrm{rev}}}{T}=0\qquad\text{（可逆なサイクル）}\]
$\delta Q_{\mathrm{rev}}$ の rev は reversible（可逆）の頭の文字です。可逆なサイクルでは、働く物の温度と熱の源の温度は同じなので、$T$ は働く物の温度です。

\noindent\begin{minipage}[c]{0.44\linewidth}\centering
\begin{tikzpicture}[x=1mm, y=1mm]
  \draw[bk arrow] (0,0) -- (48,0) node[right, font=\sffamily\footnotesize] {$V$};
  \draw[bk arrow] (0,0) -- (0,34) node[above, font=\sffamily\footnotesize] {$p$};
  \fill[bkInk] (8,26) circle (0.9) (40,8) circle (0.9);
  \node[font=\sffamily\footnotesize, anchor=south east] at (8,26) {1};
  \node[font=\sffamily\footnotesize, anchor=north west] at (40,8) {2};
  \draw[bkBlue, line width=1.2pt, -{Stealth[length=2mm]}] (8,26) .. controls (24,32) and (40,22) .. (40,8.5);
  \draw[bkOrange, line width=1.2pt, -{Stealth[length=2mm]}] (8,26) .. controls (10,12) and (24,6) .. (39.5,8);
  \node[font=\sffamily\footnotesize, text=bkBlue!80!black] at (33,27) {道A};
  \node[font=\sffamily\footnotesize, text=bkOrange!75!black] at (15,9) {道B};
\end{tikzpicture}
\end{minipage}\hfill
\begin{minipage}[c]{0.53\linewidth}\raggedright
状態1から状態2へ可逆な道Aで行き、可逆な道Bを逆にたどって1へもどると、1周の積分は0です。だから、1から2までの $\int\delta Q_{\mathrm{rev}}/T$ は、道Aでも道Bでも同じです。積分の値が、\textbf{道によらず、はじめとおわりだけで決まります}。
\end{minipage}

\par\vspace{4pt}
これは、力-4 ⑥の保存力の仕事が道によらなかったのと同じ形です。そこで、位置エネルギーを決めたのと同じように、基準の状態Oから状態Aまでの $\int\delta Q_{\mathrm{rev}}/T$ を、状態Aの\textbf{エントロピー} $S(\mathrm A)$ と決めます。小さな変化で書くと
\pheq{\dd S=\frac{\delta Q_{\mathrm{rev}}}{T}}{J/K ＝ J ÷ K}
です。こうして、エントロピーは、熱-2 ②の内部エネルギー $U$ と同じ、\textbf{状態量}になりました。④で使った「同じはじめとおわりなら、どの道でも $\Delta S$ は同じ」は、このことです。

\phboxhead{積分因子：$1/T$ をかけると全微分になる}
状態量の小さな変化は全微分になり、熱 $\delta Q$ は全微分ではありませんでした（熱-2 ⑦）。それなのに、$\delta Q_{\mathrm{rev}}$ を $T$ でわると、全微分 $\dd S$ になります。理想気体で、本当にそうなっているかを確かめましょう。

熱-2 ⑦の見分け方を思い出します。$X\,\dd x+Y\,\dd y$ が、ある関数 $f$ の全微分 $\dd f$ なら、$X=\partial f/\partial x$、$Y=\partial f/\partial y$ です。偏微分の順番を入れかえても同じ（8-2⑧）なので、
\[\frac{\partial X}{\partial y}=\frac{\partial^2f}{\partial y\,\partial x}=\frac{\partial^2f}{\partial x\,\partial y}=\frac{\partial Y}{\partial x}\]
が成り立ちます（ここの $X$・$Y$ は、$x$ と $y$ の関数の名前）。

\noindent\textbf{熱そのもの}：理想気体 $n$〔mol〕が、温度と体積を少し変える可逆な変化で受けとる熱は、第一法則 $\dd U=\delta Q+\delta W$ に、$\dd U=nC_V\,\dd T$（熱-3 ④）と、される仕事 $\delta W=-p\,\dd V$（熱-2 ③）を入れて、
\[\delta Q_{\mathrm{rev}}=nC_V\,\dd T+p\,\dd V=nC_V\,\dd T+\frac{nRT}{V}\,\dd V\]
です。$\dd T$ の係数 $nC_V$ を $V$ で偏微分すると0、$\dd V$ の係数 $nRT/V$ を $T$ で偏微分すると $nR/V$ です。等しくないので、$\delta Q$ は全微分では\textbf{ありません}。熱が道で変わる（熱-2 ③）ことの、式での表れです。

\noindent\textbf{熱÷温度}：$T$ でわると
\[\frac{\delta Q_{\mathrm{rev}}}{T}=\frac{nC_V}{T}\,\dd T+\frac{nR}{V}\,\dd V\]
$\dd T$ の係数 $nC_V/T$ を $V$ で偏微分すると0、$\dd V$ の係数 $nR/V$ を $T$ で偏微分しても0。等しいので、全微分です。実際、
\[\frac{\delta Q_{\mathrm{rev}}}{T}=\dd\bigl(nC_V\ln T+nR\ln V\bigr)\]
です（$\ln$ の微分、5-14）。このように、全微分でない式にかけると全微分になる関数を、\textbf{積分因子}といいます。\bkidx{せきぶんいんし}{積分因子}$1/T$ は、熱 $\delta Q_{\mathrm{rev}}$ の積分因子です。

\begin{phmijika}[数学の式で見る積分因子]
熱-2 例題11 で調べた $-y\,\dd x+x\,\dd y$（力-4 例題13 の力 $(-y,\ x)$ と同じ形）は、$\tfrac{\partial(-y)}{\partial y}=-1$、$\tfrac{\partial x}{\partial x}=1$ で、全微分ではありませんでした。これに $1/x^2$ をかけると、$-(y/x^2)\dd x+(1/x)\dd y$ で、$\tfrac{\partial}{\partial y}\bigl(-(y/x^2)\bigr)=-1/x^2$、$\tfrac{\partial}{\partial x}(1/x)=-1/x^2$ と等しくなり、$\dd(y/x)$ になります。$1/x^2$ が積分因子です。
\end{phmijika}

\phboxhead{Python で、2つの道をたしてみる}
理想気体（2原子）\qty{1}{mol}を、\qty{300}{K}・\qty{10}{L}から\qty{600}{K}・\qty{20}{L}へ変えます。道Aは「体積一定で温めてから、\qty{600}{K}のままふくらむ」、道Bは「\qty{300}{K}のままふくらんでから、体積一定で温める」です。どちらも可逆な道を1000歩に分け、1歩ごとの熱 $\delta Q$ と、$\delta Q/T$ をたします。
\begin{phcode}
R = 8.314; Cv = 2.5 * R; N = 1000
def heat(T1, T2, V):        \# 体積一定で温める
    Q = S = 0.0
    for i in range(N):
        Ta = T1 + (T2 - T1) * i / N
        Tb = T1 + (T2 - T1) * (i + 1) / N
        dQ = Cv * (Tb - Ta)
        Q += dQ; S += dQ / ((Ta + Tb) / 2)
    return Q, S
def expand(V1, V2, T):      \# 温度一定でふくらむ
    Q = S = 0.0
    for i in range(N):
        Va = V1 + (V2 - V1) * i / N
        Vb = V1 + (V2 - V1) * (i + 1) / N
        dQ = R * T / ((Va + Vb) / 2) * (Vb - Va)
        Q += dQ; S += dQ / T
    return Q, S
a1 = heat(300, 600, 0.010); a2 = expand(0.010, 0.020, 600)
b1 = expand(0.010, 0.020, 300); b2 = heat(300, 600, 0.020)
print("A", round(a1[0] + a2[0]), round(a1[1] + a2[1], 2))
print("B", round(b1[0] + b2[0]), round(b1[1] + b2[1], 2))
\end{phcode}
\begin{phcodeout}
A 9693 20.17
B 7964 20.17
\end{phcodeout}
道ごとに、受けとる熱の合計〔J〕と、$\delta Q/T$ の和〔J/K〕が出ます。受けとる熱は道で\qty{1729}{J}も違うのに、$\delta Q/T$ の和はどちらも\qty{20.17}{J/K}で同じです。これは、積分因子の式から出る $nC_V\ln2+nR\ln2\approx\qty{20.17}{J/K}$ と一致します。

\begin{seikaku}
理想気体では、⑦の式で確かめたとおり $1/T$ が積分因子です。どんな物質でも、熱の積分因子として同じ $1/T$ が使えることは、理想気体の性質からではなく、クラウジウスの不等式（つまり第二法則）から言えることです。ここで言う $T$ は、熱-5 ⑥の熱力学温度で、理想気体温度計の温度と同じものです。
\end{seikaku}

\phboxhead{どんな変化でも：$\Delta S\geqq\int\delta Q/T$}
状態1から状態2への、本当の（不可逆かもしれない）変化を考えます。そのあと、可逆な道で2から1へもどすと、1周のサイクルになります。クラウジウスの不等式から、
\[\int_1^2\frac{\delta Q}{T_{\mathrm b}}+\int_2^1\frac{\delta Q_{\mathrm{rev}}}{T}\leqq0\qquad\text{だから}\qquad\int_1^2\frac{\delta Q}{T_{\mathrm b}}\leqq S_2-S_1\]
です。$T_{\mathrm b}$ は、熱が通る\textbf{境目の温度}です（b は boundary（境目）の頭の文字）。左辺は「熱といっしょに入ってきたエントロピー」で、右辺の「たまったエントロピー」のほうが大きいか、等しいのです。その差を、\textbf{生まれるエントロピー}（エントロピー生成）$S_{\mathrm{gen}}$ とよびます（gen は generation の頭の文字）。\bkidx{えんとろぴーせいせい}{エントロピー生成}
\pheq{\Delta S=\int\frac{\delta Q}{T_{\mathrm b}}+S_{\mathrm{gen}}\qquad S_{\mathrm{gen}}\geqq0}{J/K}

\begin{bkfigure}
\begin{tikzpicture}[x=1mm, y=1mm]
  \filldraw[fill=bkTheme!6, draw=bkTheme, line width=1pt, rounded corners=2mm] (0,0) rectangle (52,26);
  \node[font=\sffamily\small, align=center] at (26,17.5) {たまる\\$\Delta S$};
  \node[font=\sffamily\footnotesize, align=center, fill=bkOrange!15, rounded corners=1mm, inner sep=1.5pt] at (26,5.5) {生じる $S_{\mathrm{gen}}\geqq0$};
  \draw[bk arrow, line width=1.2pt, bkRed] (-22,15) -- (-1,15);
  \node[font=\sffamily\footnotesize, anchor=south, align=center] at (-11,15.5) {入る\\$Q_{\text{入}}/T_{\text{入}}$};
  \draw[bk arrow, line width=1.2pt, bkBlue] (53,15) -- (74,15);
  \node[font=\sffamily\footnotesize, anchor=south, align=center] at (63.5,15.5) {出る\\$Q_{\text{出}}/T_{\text{出}}$};
  \node[font=\sffamily\footnotesize, anchor=west, align=left] at (78,13) {入る $-$ 出る $+$ 生じる\\＝ たまる};
\end{tikzpicture}
\bkfigcaption{エントロピーのバランス式。エネルギー（熱-2）と同じ「入る－出る＋生じる＝たまる」の形だが、エントロピーは「生じる」が0以上で、なくなることはない。}\label{ph:熱-6:fig-balance}
\end{bkfigure}

\phboxhead{孤立した系：エントロピー増大の法則}
まわりと熱をやりとりしない（孤立した）系では、$\int\delta Q/T_{\mathrm b}=0$ なので、$\Delta S=S_{\mathrm{gen}}\geqq0$ です。②で予想したエントロピー増大の法則が、熱-5 の第二法則から導けました。系とまわりを合わせて1つの孤立した系と見ると、
\[S_{\mathrm{gen}}=\Delta S_{\text{系}}+\Delta S_{\text{まわり}}\geqq0\]
です。\textbf{系だけなら、エントロピーは減ってもかまいません}（冷蔵庫の中の食べ物、凍る水）。でも、そのときは、まわりのエントロピーが、それ以上に増えています。

\begin{phflow}[生まれるエントロピー]
\phflowrow{$S_{\mathrm{gen}}=\Delta S_{\text{系}}+\Delta S_{\text{まわり}}\geqq0$}{可逆な変化\\（摩擦・温度差なし）}{$S_{\mathrm{gen}}=0$、$\oint\dfrac{\delta Q}{T}=0$}
\phflowrow{$\Delta S=\displaystyle\int\frac{\delta Q}{T_{\mathrm b}}+S_{\mathrm{gen}}$}{孤立した系\\（熱の出入りなし）}{$\Delta S\geqq0$}
\end{phflow}

%% ---------------------------------------------------------------
\bksubsection{7}{水・氷・気体のエントロピーの式}
\label{ph:熱-6:sec7}
$\dd S=\delta Q_{\mathrm{rev}}/T$ を、よく出てくる場合に積分します。どれも「可逆に熱を受けとる道」を1つ考えて計算します。

\phboxhead{温度一定のとき：$\Delta S=Q/T$}
$T$ が一定なら、$T$ を積分の外に出せて、$\Delta S=(1/T)\int\delta Q_{\mathrm{rev}}=Q/T$ です。②で決めた式が、一般の式の特別な場合として出てきました。

\phboxhead{水や金属を温める：$\Delta S=mc\ln(T_2/T_1)$}
液体や固体は、温めても体積がほとんど変わらないので、受けとる熱は $\delta Q_{\mathrm{rev}}=mc\,\dd T$ です（熱-1 ③）。比熱 $c$ が一定なら、
\pheq{\Delta S=\int_{T_1}^{T_2}\frac{mc}{T}\,\dd T=mc\Bigl[\ln T\Bigr]_{T_1}^{T_2}=mc\ln\frac{T_2}{T_1}}{kg × J/(kg·K) ＝ J/K}
です（5-14 の $\int(1/x)\dd x$）。③の表計算は、この積分を、細い長方形（\qty{1}{K}ずつ）の和で数えたものでした。

\begin{phyomi}[$\Delta S=mc\ln\dfrac{T_2}{T_1}$]
「温度の\textbf{差}ではなく、\textbf{比}で決まる」。比熱が同じなら、\qty{280}{K}から\qty{300}{K}に温めても、\qty{350}{K}から\qty{375}{K}に温めても（どちらも $15/14$ 倍）、同じだけ増える。低い温度での\qty{1}{K}は、比で見ると大きな変化だから。
\end{phyomi}

\begin{phbaisuru}{質量 $m$}
同じ温度の変化なら、$\Delta S$ は2倍です。エントロピーは、内部エネルギーと同じく、量に比例する量です。$T_2$ を変えると、$\Delta S$ は $\ln T_2$ にそってゆっくり増えます。
\par\smallskip\noindent\begin{tikzpicture}
\begin{axis}[bkplot, width=100mm, height=44mm, xmin=0, xmax=82, ymin=0, ymax=2300, xtick={20,40,60,80}, ytick={500,1000,2000}, xlabel={温めたあとの温度（℃）}, ylabel={$\Delta S$（J/K）}]
  \addplot[bkcurve, domain=0:80] {4200*ln((x+273.15)/273.15)};
  \addplot[bkcurve2, domain=0:80] {8400*ln((x+273.15)/273.15)};
  \node[font=\sffamily\footnotesize, anchor=west, text=bkBlue!80!black] at (axis cs:60,700) {\qty{1}{kg}};
  \node[font=\sffamily\footnotesize, anchor=west, text=bkOrange!75!black] at (axis cs:20,1500) {\qty{2}{kg}（2倍）};
  \node[font=\sffamily\footnotesize, anchor=west, text=bkGray] at (axis cs:46,250) {0\,℃の水を温める};
\end{axis}
\end{tikzpicture}
\end{phbaisuru}

\phboxhead{とける・わく：$\Delta S=mL/T$}
氷がとける間、水がわく間は、温度が一定のまま熱を受けとります（熱-1 ⑤）。潜熱を $L$〔J/kg〕とすると、$Q=mL$ なので、
\[\Delta S=\frac{mL}{T}\]
です。氷\qty{1}{kg}がとけると $334000\div273.15\approx\qty{1.22e3}{J/K}$（例題2）、水\qty{1}{kg}が100\,℃でわくと $2257000\div373.15\approx\qty{6.05e3}{J/K}$ も増えます。分子が自由に動けるようになるほど、エントロピーは大きく増えます（熱-15）。

\phboxhead{理想気体}
⑥の積分因子の式をそのまま積分すると、
\pheq{\Delta S=nC_V\ln\frac{T_2}{T_1}+nR\ln\frac{V_2}{V_1}}{J/K}
です。温度が一定なら2つめの項だけが残り、④の $nR\ln2$ になります。1 kg あたりの\textbf{比エントロピー} $s$ では、$nC_V$ を $c_v$、$nR$ を $R_{\mathrm s}$（熱-3 ③）に変えます。さらに、$pV=nRT$ から $V_2/V_1=(T_2/T_1)\cdot(p_1/p_2)$ を入れ、$c_p=c_v+R_{\mathrm s}$（熱-3 ④）を使うと、圧力で書いた形
\pheq{\Delta s=c_p\ln\frac{T_2}{T_1}-R_{\mathrm s}\ln\frac{p_2}{p_1}}{J/(kg·K)}
になります。圧縮機やタービンでは、温度と圧力がわかるので、こちらの形をよく使います（例題12）。

\phboxhead{断熱で可逆なら：$s$ が一定}
熱の出入りがなく（$\delta Q=0$）、可逆なら、$\dd S=0$ で、エントロピーは変わりません。これを\textbf{等エントロピー変化}といいます。理想気体の式で $\Delta s=0$ とおくと、
\[c_v\ln\frac{T_2}{T_1}=-R_{\mathrm s}\ln\frac{V_2}{V_1}\qquad\text{だから}\qquad TV^{R_{\mathrm s}/c_v}=\text{一定}\]
$R_{\mathrm s}/c_v=(c_p-c_v)/c_v=\gamma-1$ なので $TV^{\gamma-1}$ が一定、$pV=nRT$ を使うと $pV^{\gamma}$ が一定です。熱-3 ⑥で変数分離から導いたポアソンの式に、別の道から着きました（例題10）。

\begin{phflow}[エントロピーの式]
\phflowrow{$\dd S=\dfrac{\delta Q_{\mathrm{rev}}}{T}$}{温度 $T$ が一定}{$\Delta S=\dfrac{Q}{T}$}
\phflowrow{$\dd S=\dfrac{mc\,\dd T}{T}$}{液体・固体\\$c$ が一定}{$\Delta S=mc\ln\dfrac{T_2}{T_1}$}
\phflowrow{$\dd s=\dfrac{c_v\,\dd T}{T}+\dfrac{R_{\mathrm s}\,\dd V}{V}$}{断熱で可逆\\（$\dd s=0$）}{$pV^{\gamma}=$ 一定}
\end{phflow}

%% ---------------------------------------------------------------
\bksubsection{8}{$T$-$s$ 線図：面積は熱、閉じた道の面積は仕事}
\label{ph:熱-6:sec8}
\phboxhead{エントロピーを横に、温度をたてにとる}
$p$-$V$ 図（熱-2）では、道の下の面積 $\int p\,\dd V$ が、気体がする仕事でした。同じように、エントロピー $S$ を横、温度 $T$ をたてにとった図を\textbf{$T$-$S$ 線図}（1 kg あたりなら $T$-$s$ 線図）といいます。\bkidx{てぃーえすせんず}{$T$-$s$ 線図}$\dd S=\delta Q_{\mathrm{rev}}/T$ から $\delta Q_{\mathrm{rev}}=T\,\dd S$ なので、可逆な道の下の面積は
\[\int T\,\dd S=Q_{\mathrm{rev}}\qquad\text{（受けとる熱）}\]
です。右へ進む（$S$ が増える）ときは熱を受けとり、左へ進むときは熱を捨てます。

\begin{bkfigure}
\begin{tikzpicture}[x=1mm, y=1mm]
  \begin{scope}
  \draw[bk arrow] (0,0) -- (54,0) node[right, font=\sffamily\footnotesize] {$S$};
  \draw[bk arrow] (0,0) -- (0,40) node[above, font=\sffamily\footnotesize] {$T$};
  \fill[bkRed!15] (10,0) -- (10,28) -- (40,28) -- (40,0) -- cycle;
  \draw[bkRed, line width=1.4pt, -{Stealth[length=2.2mm]}] (10,28) -- (40,28);
  \draw[bkBlue, line width=1.4pt, -{Stealth[length=2.2mm]}] (46,30) -- (46,8);
  \node[font=\sffamily\footnotesize, anchor=south] at (25,28.5) {等温（横）};
  \node[font=\sffamily\footnotesize, align=center, anchor=west] at (46.5,20) {断熱で\\可逆\\（たて）};
  \node[font=\sffamily\footnotesize, align=center] at (25,13) {面積＝\\受けとる熱};
  \node[font=\sffamily\bfseries\footnotesize, anchor=south] at (28,41) {（ア）線の向きと面積};
  \end{scope}
  \begin{scope}[shift={(72,0)}]
  \draw[bk arrow] (0,0) -- (54,0) node[right, font=\sffamily\footnotesize] {$S$};
  \draw[bk arrow] (0,0) -- (0,40) node[above, font=\sffamily\footnotesize] {$T$};
  \fill[bkOrange!20] (10,12) rectangle (40,30);
  \fill[bkBlue!10] (10,0) rectangle (40,12);
  \draw[bkInk, line width=1.3pt, -{Stealth[length=2mm]}] (10,30) -- (40,30);
  \draw[bkInk, line width=1.3pt, -{Stealth[length=2mm]}] (40,30) -- (40,12);
  \draw[bkInk, line width=1.3pt, -{Stealth[length=2mm]}] (40,12) -- (10,12);
  \draw[bkInk, line width=1.3pt, -{Stealth[length=2mm]}] (10,12) -- (10,30);
  \node[font=\sffamily\footnotesize, anchor=east] at (-1,30) {$T_{\mathrm H}$};
  \node[font=\sffamily\footnotesize, anchor=east] at (-1,12) {$T_{\mathrm L}$};
  \draw[bkGray, dashed, line width=0.4pt] (0,30) -- (10,30) (0,12) -- (10,12);
  \node[font=\sffamily\footnotesize, align=center] at (25,21) {正味の仕事\\$W_{\text{出}}$};
  \node[font=\sffamily\footnotesize] at (25,6) {捨てる熱 $Q_{\mathrm L}$};
  \node[font=\sffamily\bfseries\footnotesize, anchor=south] at (28,41) {（イ）カルノーサイクル};
  \end{scope}
\end{tikzpicture}
\bkfigcaption{$T$-$S$ 線図。（ア）等温の変化は横の線、断熱で可逆な変化はたての線。可逆な道の下の面積は受けとる熱。（イ）カルノーサイクルは長方形で、囲んだ面積が正味の仕事、その下の面積が捨てる熱。}\label{ph:熱-6:fig-ts}
\end{bkfigure}

\phboxhead{カルノーサイクルは長方形}
カルノーサイクルは、2つの等温と、2つの断熱で可逆な変化でできているので、$T$-$S$ 線図では長方形です。上の辺の下の面積 $T_{\mathrm H}\Delta S$ が受けとる熱 $Q_{\mathrm H}$、下の辺の下の面積 $T_{\mathrm L}\Delta S$ が捨てる熱 $Q_{\mathrm L}$、その差の長方形が正味の仕事です。効率は
\[\eta=\frac{(T_{\mathrm H}-T_{\mathrm L})\Delta S}{T_{\mathrm H}\Delta S}=1-\frac{T_{\mathrm L}}{T_{\mathrm H}}\]
で、熱-5 ⑤のカルノーの効率が、図の長方形の面積の比から1行で出てきます（例題6）。

\phboxhead{閉じた道の面積は、どちらの図でも正味の仕事}
\label{ph:熱-6:menseki}数学の 8-7 のコラムでは、エンジンの気体が1回にする仕事は $p$-$V$ 図の閉じた道の線積分 $\oint p\,\dd V$ で、道が囲む面積だと学びました（そこには「物理 熱-6」と書いてありました）。$T$-$S$ 線図でも、閉じた道が囲む面積 $\oint T\,\dd S$ は、1周で受けとる正味の熱です。1周すると内部エネルギーはもとにもどる（$\oint\dd U=0$、熱-2 ⑧）ので、第一法則から、
\pheq{\oint T\,\dd S=\oint p\,\dd V=W_{\text{出}}\ \text{（1周で外にする正味の仕事）}}{J}
です。\textbf{$p$-$V$ 図で囲んだ面積と、$T$-$S$ 線図で囲んだ面積は、同じ大きさ}で、どちらも正味の仕事です。道を時計回りに回るとき、外へ仕事をする熱機関、反時計回りのとき、仕事を入れて熱をくみ上げる冷凍機です。

\begin{bkfigure}
\begin{tikzpicture}
\begin{axis}[bkplot, name=pv, width=62mm, height=52mm, xmin=0, xmax=3.6, ymin=0, ymax=250, xtick={1,2,3}, ytick={100,200}, xlabel={$V$（L）}, ylabel={$p$（kPa）}]
  \fill[bkOrange!22] (axis cs:1,100) rectangle (axis cs:3,200);
  \draw[bkInk, line width=1.3pt, -{Stealth[length=2mm]}] (axis cs:1,200) -- (axis cs:3,200);
  \draw[bkInk, line width=1.3pt, -{Stealth[length=2mm]}] (axis cs:3,200) -- (axis cs:3,100);
  \draw[bkInk, line width=1.3pt, -{Stealth[length=2mm]}] (axis cs:3,100) -- (axis cs:1,100);
  \draw[bkInk, line width=1.3pt, -{Stealth[length=2mm]}] (axis cs:1,100) -- (axis cs:1,200);
  \node[font=\sffamily\footnotesize] at (axis cs:2,150) {\qty{200}{J}};
  \node[font=\sffamily\footnotesize, anchor=south east] at (axis cs:1,201) {A};
  \node[font=\sffamily\footnotesize, anchor=south west] at (axis cs:3,201) {B};
  \node[font=\sffamily\footnotesize, anchor=north west] at (axis cs:3,99) {C};
  \node[font=\sffamily\footnotesize, anchor=north east] at (axis cs:1,99) {D};
\end{axis}
\begin{axis}[bkplot, at={(pv.east)}, anchor=west, xshift=16mm, width=62mm, height=52mm, xmin=-1.5, xmax=3.2, ymin=0, ymax=1000, xtick={-1,0,1,2,3}, ytick={150,300,450,900}, xlabel={$S$（J/K）}, ylabel={$T$（K）}, axis y line=left]
  \fill[bkOrange!30] (axis cs:0.000,300.0) -- (axis cs:0.425,360.0) -- (axis cs:0.785,420.0) -- (axis cs:1.097,480.0) -- (axis cs:1.372,540.0) -- (axis cs:1.617,600.0) -- (axis cs:1.840,660.0) -- (axis cs:2.043,720.0) -- (axis cs:2.230,780.0) -- (axis cs:2.403,840.0) -- (axis cs:2.564,900.0) -- (axis cs:2.478,855.0) -- (axis cs:2.388,810.0) -- (axis cs:2.293,765.0) -- (axis cs:2.191,720.0) -- (axis cs:2.084,675.0) -- (axis cs:1.969,630.0) -- (axis cs:1.845,585.0) -- (axis cs:1.712,540.0) -- (axis cs:1.567,495.0) -- (axis cs:1.408,450.0) -- (axis cs:1.247,420.0) -- (axis cs:1.074,390.0) -- (axis cs:0.887,360.0) -- (axis cs:0.684,330.0) -- (axis cs:0.462,300.0) -- (axis cs:0.216,270.0) -- (axis cs:-0.059,240.0) -- (axis cs:-0.370,210.0) -- (axis cs:-0.730,180.0) -- (axis cs:-1.155,150.0) -- (axis cs:-0.996,165.0) -- (axis cs:-0.851,180.0) -- (axis cs:-0.718,195.0) -- (axis cs:-0.594,210.0) -- (axis cs:-0.479,225.0) -- (axis cs:-0.372,240.0) -- (axis cs:-0.271,255.0) -- (axis cs:-0.175,270.0) -- (axis cs:-0.085,285.0) -- cycle;
  \addplot[bkInk, line width=1.3pt, domain=300:900, samples=40] ({2.3335*ln(x/300)},{x});
  \addplot[bkInk, line width=1.3pt, domain=450:900, samples=40] ({2.5634+1.6668*ln(x/900)},{x});
  \addplot[bkInk, line width=1.3pt, domain=150:450, samples=40] ({1.4082+2.3335*ln(x/450)},{x});
  \addplot[bkInk, line width=1.3pt, domain=150:300, samples=40] ({-1.1552+1.6668*ln(x/150)},{x});
  \fill[bkInk] (axis cs:0,300) circle (1.6pt) (axis cs:2.5634,900) circle (1.6pt) (axis cs:1.4082,450) circle (1.6pt) (axis cs:-1.1552,150) circle (1.6pt);
  \node[font=\sffamily\footnotesize, anchor=south east] at (axis cs:-0.05,305) {A};
  \node[font=\sffamily\footnotesize, anchor=south east] at (axis cs:2.5,905) {B};
  \node[font=\sffamily\footnotesize, anchor=north west] at (axis cs:1.45,445) {C};
  \node[font=\sffamily\footnotesize, anchor=north west] at (axis cs:-1.1,145) {D};
  \draw[bkInk, line width=0.5pt] (axis cs:1.0,420) -- (axis cs:2.0,250);
  \node[font=\sffamily\footnotesize, anchor=north west, inner sep=1pt] at (axis cs:1.9,250) {\qty{200}{J}};
\end{axis}
\end{tikzpicture}
\bkfigcaption{熱-2 例題12 の長方形のサイクル（2原子の理想気体 \qty{0.080}{mol}、Aで \qty{300}{K}）を、$p$-$V$ 図（左）と $T$-$S$ 線図（右。Aを $S=0$ とした）にかいた。形は違うが、囲んだ面積はどちらも \qty{200}{J}。}\label{ph:熱-6:fig-pvts}
\end{bkfigure}

右の図の4本の線は、⑦の式から計算しました。定圧の A→B では $\Delta S=nC_p\ln(T/T_{\mathrm A})$、定積の B→C では $\Delta S=nC_V\ln(T/T_{\mathrm B})$ です。A→B で受けとる熱は\qty{1400}{J}、D→A で\qty{250}{J}、B→C で捨てる熱は\qty{750}{J}、C→D で\qty{700}{J}で、正味の熱は $1400+250-750-700=\qty{200}{J}$。$p$-$V$ 図の長方形の面積 $\qty{100}{kPa}\times\qty{2}{L}=\qty{200}{J}$ と同じです。

\begin{phyomi}[$\oint T\,\dd S=\oint p\,\dd V$]
「1周で入った正味の熱は、1周で出した正味の仕事」。$p$-$V$ 図は仕事が、$T$-$S$ 線図は熱が、面積で見える図。サイクルを考えるときは、2つを並べてかくとよい（熱-7・熱-8）。
\end{phyomi}

%% ---------------------------------------------------------------
\bksubsection{9}{エントロピー生成とエクセルギー：損失を数える}
\label{ph:熱-6:sec9}
\phboxhead{熱のうち、仕事にできる分：エクセルギー}
温度 $T$ の熱 $Q$ から、いちばん多く取り出せる仕事はいくらでしょう。まわり（大気や海）の温度を $T_0$ とすると、$T$ と $T_0$ の間で可逆な機関を動かすときが最大で、
\pheq{E_{\mathrm x}=Q\Bigl(1-\frac{T_0}{T}\Bigr)}{J}
です（熱-5 ⑤のカルノーの効率）。これを熱の\textbf{エクセルギー}（有効エネルギー）といいます。\bkidx{えくせるぎー}{エクセルギー}残りの $Q\,T_0/T$ は、どうしてもまわりへ捨てる熱です。エネルギーは保存しますが、エクセルギーは不可逆な変化で減っていきます。

\begin{phbaisuru}{熱の温度 $T$（まわりは $T_0=\qty{300}{K}$）}
$T$ を\qty{600}{K}から\qty{1200}{K}にすると、エクセルギーの割合 $1-T_0/T$ は0.50から0.75に上がります。2倍にはなりません。$T$ が $T_0$ に近い熱（ぬるい熱）は、どんなに量が多くても、仕事にできる分はわずかです。
\par\smallskip\noindent\begin{tikzpicture}
\begin{axis}[bkplot, width=100mm, height=42mm, xmin=250, xmax=1300, ymin=0, ymax=1, xtick={300,600,900,1200}, ytick={0.5,0.75,1}, xlabel={$T$（K）}, ylabel={$1-T_0/T$}]
  \addplot[bkcurve, domain=300:1280] {1-300/x};
  \fill[bkRed] (axis cs:600,0.5) circle (1.8pt) (axis cs:1200,0.75) circle (1.8pt);
  \node[font=\sffamily\footnotesize, anchor=north west, text=bkRed!75!black] at (axis cs:620,0.48) {600 K で 0.50};
  \node[font=\sffamily\footnotesize, anchor=north west, text=bkRed!75!black] at (axis cs:1000,0.70) {1200 K で 0.75};
\end{axis}
\end{tikzpicture}
\end{phbaisuru}

\phboxhead{失われる仕事＝$T_0\times$生まれるエントロピー}
温度 $T_{\mathrm H}$ の熱 $Q$ が、仕事をせずに、温度 $T_{\mathrm L}$ の所へ流れたとします（②）。流れる前のエクセルギーは $Q(1-T_0/T_{\mathrm H})$、流れたあとは $Q(1-T_0/T_{\mathrm L})$ なので、失われたエクセルギーは
\[Q\Bigl(1-\frac{T_0}{T_{\mathrm H}}\Bigr)-Q\Bigl(1-\frac{T_0}{T_{\mathrm L}}\Bigr)=T_0\,Q\Bigl(\frac{1}{T_{\mathrm L}}-\frac{1}{T_{\mathrm H}}\Bigr)=T_0\,S_{\mathrm{gen}}\]
です。かっこの中は、②で求めた、熱が流れて生まれたエントロピーそのものです。一般に、
\pheq{\text{失われる仕事}=T_0\,S_{\mathrm{gen}}}{J ＝ K × J/K}
が成り立ちます（グイ・ストドラの定理）。\textbf{生まれるエントロピーに、まわりの温度をかけると、取り出しそこねた仕事になる}のです。摩擦・混ぜる・温度差のある熱の流れ・自由膨張など、エントロピーが生まれる所が、機械の「むだ」の正体です。

\begin{bkfigure}
\begin{tikzpicture}[x=1mm, y=1mm]
  \node[font=\sffamily\bfseries\footnotesize, anchor=west] at (0,30) {電気ストーブ（\qty{1}{kW}）};
  \fill[bkOrange!70] (0,20) rectangle (6.8,26);
  \fill[bkGray!35] (6.8,20) rectangle (100,26);
  \draw[bkInk, line width=0.5pt] (0,20) rectangle (100,26);
  \node[font=\sffamily\footnotesize, anchor=north west] at (0,19.3) {部屋を暖める分 \qty{68}{W}};
  \node[font=\sffamily\footnotesize, anchor=north east] at (100,19.3) {失われるエクセルギー \qty{932}{W}};
  \node[font=\sffamily\bfseries\footnotesize, anchor=west] at (0,8) {ヒートポンプ（COP 4、電気 \qty{250}{W}）};
  \fill[bkOrange!70] (0,-2) rectangle (6.8,4);
  \fill[bkGray!35] (6.8,-2) rectangle (25,4);
  \draw[bkInk, line width=0.5pt] (0,-2) rectangle (25,4);
  \node[font=\sffamily\footnotesize, anchor=west] at (27,1) {\qty{68}{W}は同じ。失われる分は \qty{182}{W} だけ};
\end{tikzpicture}
\bkfigcaption{20\,℃の部屋に \qty{1}{kW} の熱を入れるとき（外は 0\,℃）の、エクセルギーの帯。帯の長さは入れた電気（どちらもエクセルギーはそのまま）。部屋の熱 \qty{1}{kW} のエクセルギーは $1000\times(1-273.15/293.15)\approx\qty{68}{W}$。熱-0 の熱の流れの帯（エネルギー）では、どちらも「\qty{1}{kW}の熱」で同じに見える。}\label{ph:熱-6:fig-exergy}
\end{bkfigure}

\begin{phsaki}{ゆっくりでない変化の熱力学}
この章の熱力学は、はじめとおわりが「止まって落ち着いた状態（平衡状態）」の変化を扱いました。熱が流れ続ける板、電流が流れ続ける導線のように、ずっと流れ続ける状態では、1秒あたりに生まれるエントロピーを、流れの大きさと、流れを起こす原因（温度の差など）の積の和で書きます。オンサーガーは1931年、2種類の流れが互いに影響し合うときの係数の間に、決まった関係（相反定理）があることを示しました（1968年にノーベル化学賞）。熱電素子（熱-5）のしくみもこの考えで説明できます。このような\textbf{非平衡熱力学}は、専門の本で扱います。エントロピーを「並び方の数」で数える見方は、熱-15 で学びます。
\end{phsaki}

\begin{motto}[5]
\phboxhead{衝撃波でもエントロピーが増える（熱-4 ⑨）}
超音速の流れの中には、圧力・温度・密度が、ごく薄い所で急にはね上がる\textbf{衝撃波}ができます（熱-4 ⑨）。その薄い所では、気体の粘性と熱伝導が強くはたらき、エントロピーが生まれます。質量・運動量・エネルギーのバランス式（力-13・熱-4）だけを見ると、「圧力が急に下がる衝撃波」も式の上ではありえますが、計算するとエントロピーが減ってしまいます。だから、そのような衝撃波は起きず、自然にできるのは圧力が上がる衝撃波だけです。第二法則が、どちらの解が本物かを選んでいる例です。
\end{motto}

\begin{senmonsho}[熱力学の教科書ふうの書き方]
\begin{bkhonbun}
閉じた系が温度 $T$ の熱源から熱 $d'Q$ を受け取るとき，$dS\geqq d'Q/T$ が成り立つ（クラウジウスの不等式）．等号は準静的過程に限る．孤立系では $dS\geqq0$ である．周囲温度 $T_0$ のもとでの損失仕事は $W_{\mathrm{lost}}=T_0S_{\mathrm{gen}}$ で与えられる．
\end{bkhonbun}
\bkyaku{$d'Q$（ディー・ダッシュ・キュー）は、この本の $\delta Q$ と同じ「全微分でない小さな熱」。「熱源」は大きな水そうのように温度が変わらない熱の源。「準静的過程」は、ゆっくり進む変化で、この章では可逆な変化と同じに使ってよい。「閉じた系」は物の出入りのない系（熱と仕事は出入りしてよい）、「孤立系」は熱も仕事も出入りしない系。lost は「失われた」。}
\end{senmonsho}

\begin{taisetsu}
\def\tsl#1{\makebox[7.4\zw][l]{\sffamily\bfseries\small\color{\bkTT}#1}}
\tsl{エントロピー}$\dd S=\dfrac{\delta Q_{\mathrm{rev}}}{T}$。温度一定なら $\Delta S=\dfrac{Q}{T}$\\[3pt]
\tsl{第二法則}$\oint\dfrac{\delta Q}{T}\leqq0$。$S_{\mathrm{gen}}=\Delta S_{\text{系}}+\Delta S_{\text{まわり}}\geqq0$\\[3pt]
\tsl{よく使う式}$mc\ln\dfrac{T_2}{T_1}$、$\dfrac{mL}{T}$、$c_p\ln\dfrac{T_2}{T_1}-R_{\mathrm s}\ln\dfrac{p_2}{p_1}$\\[3pt]
\tsl{$T$-$S$ 線図}道の下の面積＝熱。囲んだ面積＝正味の仕事（$p$-$V$ 図と同じ）\\[3pt]
\tsl{エクセルギー}$Q\Bigl(1-\dfrac{T_0}{T}\Bigr)$。失われる仕事＝$T_0S_{\mathrm{gen}}$
\end{taisetsu}

\begin{kotoba}
\small
\bkword{可逆変化}{まわりに何も変化を残さずに、もとにもどせる変化。}
\bkword{エントロピー}{可逆に受けとった熱を温度でわって足した量。状態量。単位 J/K。}
\bkword{クラウジウスの不等式}{1周の $\oint\delta Q/T$ は0以下。可逆なら0。}
\bkword{積分因子}{全微分でない式にかけて、全微分にする関数。熱には $1/T$。}
\bkword{エントロピー生成}{不可逆な変化で生まれるエントロピー。いつも0以上。}
\bkword{エクセルギー}{まわりの温度 $T_0$ のもとで、取り出せるいちばん大きな仕事。}
\end{kotoba}

\bkkigocards{\bkkigocard{$S$}{エス}{エントロピー〔J/K〕。可逆に受けとった熱÷絶対温度だけ増える。1 kg あたりは $s$〔J/(kg·K)〕}{entropy}{熱-6}}
{\bkkigocard{$S_{\mathrm{gen}}$}{エス・ジェン}{生まれるエントロピー。いつも0以上。$T_0S_{\mathrm{gen}}$ が失われる仕事}{entropy generation}{熱-6}}

%% ======================== 例題 ========================
\Needspace{10\baselineskip}\bksection{reidai}{例題でたしかめよう}
\label{ph:熱-6:ex1}
\begin{phreidai}[お湯と水を混ぜたときのエントロピー]
80\,℃のお湯\qty{1.0}{kg}と、20\,℃の水\qty{1.0}{kg}を混ぜます。混ぜたあとの温度と、全体のエントロピーの変化を求めましょう（\bkhatenaref）。
\phbamen
温度の違う水が混ざって、温度がそろう、もどらない変化です。エネルギーは保存し、エントロピーは増えるはずです。
\phmushi
なべとまわりの空気へ逃げる熱と、混ぜる仕事を無視します。比熱 $c$ は一定で、\qty{4.2}{kJ/(kg.K)}とします。
\phzu
お湯：$m=\qty{1.0}{kg}$、$T_1=\qty{353.15}{K}$。水：$m=\qty{1.0}{kg}$、$T_2=\qty{293.15}{K}$。混ぜたあと：$T_{\mathrm f}$。
\phriyu
温度は、熱のバランス式（熱-1 ④）で決まります。エントロピーは状態量なので、それぞれの水を、ゆっくり温度 $T_{\mathrm f}$ まで変える可逆な道で数えてよく、⑦の $mc\ln(T_2/T_1)$ が使えます。
\phtateru
\[mc(T_{\mathrm f}-T_1)+mc(T_{\mathrm f}-T_2)=0,\qquad \Delta S=mc\ln\frac{T_{\mathrm f}}{T_1}+mc\ln\frac{T_{\mathrm f}}{T_2}\]
\phtoku
同じ量なので $T_{\mathrm f}=(T_1+T_2)/2=\qty{323.15}{K}$（50\,℃）です。
\[\Delta S=4200\ln\frac{323.15}{353.15}+4200\ln\frac{323.15}{293.15}\approx-373+409=\qty{36}{J/K}\]
\phtashikame
単位：J/(kg·K)×kg＝J/K。\\
符号：お湯は減り、水は増え、合計は正。③の表計算の\qty{36.35}{J/K}と合っています。\\
極端な場合：2つが同じ温度なら、$T_{\mathrm f}=T_1=T_2$ で $\Delta S=0$（何も起きない）。
\kotae 50\,℃、全体のエントロピーは約\qty{36}{J/K}増える
\end{phreidai}
\begin{jibun}
60\,℃のお湯\qty{1.0}{kg}と、20\,℃の水\qty{1.0}{kg}を混ぜたときの、混ぜたあとの温度と、全体のエントロピーの変化を求めましょう（$c=\qty{4.2}{kJ/(kg.K)}$）。
\bkjans{40\,℃。お湯 $4200\ln(313.15/333.15)\approx-260$、水 $4200\ln(313.15/293.15)\approx277$ で、合計 約\qty{17}{J/K}増える}
\end{jibun}

\label{ph:熱-6:ex2}
\begin{phreidai}[氷がとけるときのエントロピー]
0\,℃の氷\qty{1.0}{kg}が、とけて0\,℃の水になるとき、氷（水）のエントロピーの変化を求めましょう。融解熱は\qty{334}{kJ/kg}とします。
\phbamen
氷水の中で、氷がとけていく場面です。とける間、温度は0\,℃のままです（熱-1 ⑤）。
\phmushi
とける間に、まわりとの温度の差がとても小さく、ゆっくり熱を受けとる（可逆）とします。圧力は1気圧で一定です。
\phzu
$m=\qty{1.0}{kg}$、$L=\qty{334}{kJ/kg}$、$T=\qty{273.15}{K}$。
\phriyu
温度が一定のまま熱を受けとるので、⑦の $\Delta S=mL/T$（②の $Q/T$）が使えます。
\phtateru
\[\Delta S=\frac{mL}{T}\]
\phtoku
\[\Delta S=\frac{1.0\times334000}{273.15}\approx\qty{1.22e3}{J/K}\]
\phtashikame
単位：J÷K＝J/K。\\
大きさ：例題1の約34倍です。とけると分子が自由に動けるようになり、エントロピーが大きく増えます。\\
逆に、0\,℃の水が凍るときは、同じ大きさだけ減ります（章末問題2）。
\kotae 約\qty{1.22e3}{J/K}増える
\end{phreidai}
\begin{jibun}
0\,℃の氷\qty{0.20}{kg}が、とけて0\,℃の水になるときの、エントロピーの変化を求めましょう（融解熱\qty{334}{kJ/kg}）。
\bkjans{$0.20\times334000\div273.15\approx\qty{245}{J/K}$ 増える}
\end{jibun}

\label{ph:熱-6:ex3}
\begin{phreidai}[熱が流れて生まれるエントロピー]
\qty{400}{K}の大きな物から、\qty{300}{K}の大きな物へ、熱\qty{1000}{J}が流れました。生まれたエントロピーを求めましょう。
\phbamen
熱い物と冷たい物の間を、熱が流れる場面です（図\ref{ph:熱-6:fig-rei}の（ウ））。
\phmushi
どちらの物も大きく、\qty{1000}{J}で温度はほとんど変わらないとします。まわりとの熱のやりとりは無視します。
\phzu
$T_{\mathrm H}=\qty{400}{K}$、$T_{\mathrm L}=\qty{300}{K}$、$Q=\qty{1000}{J}$。
\phriyu
それぞれの物の中では温度が一定なので、②の $\Delta S=Q/T$ が使えます。2つを合わせた全体は孤立しているので、合計の増えが生まれたエントロピーです（⑥）。
\phtateru
\[S_{\mathrm{gen}}=-\frac{Q}{T_{\mathrm H}}+\frac{Q}{T_{\mathrm L}}\]
\phtoku
\[S_{\mathrm{gen}}=-\frac{1000}{400}+\frac{1000}{300}=-2.50+3.33\approx\qty{0.83}{J/K}\]
\phtashikame
符号：正なので、この向きの流れは起きてよい。\\
極端な場合：$T_{\mathrm H}=T_{\mathrm L}$ なら0（温度の差がないと熱は流れず、生まれない）。\\
大きさ：温度の差が大きいほど、多く生まれます。
\kotae 約\qty{0.83}{J/K}
\end{phreidai}
\begin{jibun}
\qty{350}{K}の大きな物から\qty{280}{K}の大きな物へ、熱\qty{500}{J}が流れたときの、生まれたエントロピーを求めましょう。
\bkjans{$-\tfrac{500}{350}+\tfrac{500}{280}\approx-1.43+1.79$ で、約\qty{0.36}{J/K}}
\end{jibun}

\label{ph:熱-6:ex4}
\begin{phreidai}[自由膨張と、等温の可逆な膨張]
理想気体\qty{1}{mol}が、温度\qty{300}{K}のまま、体積が2倍になります。（1）等温の可逆な膨張と、（2）真空への自由膨張のそれぞれで、気体・まわり・合計のエントロピーの変化を求めましょう。
\phbamen
同じはじめとおわりの状態を、2つの違う道で結ぶ場面です（図\ref{ph:熱-6:fig-jiyuu}）。
\phmushi
気体は理想気体とします。容器の熱容量は無視します。（1）のまわりは、温度\qty{300}{K}の大きな水そうです。
\phzu
$n=\qty{1}{mol}$、$T=\qty{300}{K}$、$V_2/V_1=2$、$R=\qty{8.314}{J/(mol.K)}$。
\phriyu
気体のエントロピーは状態量なので、どちらの道でも同じで、④の $nR\ln(V_2/V_1)$ です。まわりは、熱を出した分だけ $-Q/T$ 変わります。
\phtateru
\[\Delta S_{\text{気体}}=nR\ln2,\qquad\text{（1）}\ \Delta S_{\text{まわり}}=-\frac{nRT\ln2}{T},\qquad\text{（2）}\ \Delta S_{\text{まわり}}=0\]
\phtoku
$\Delta S_{\text{気体}}=1\times8.314\times0.6931\approx\qty{5.76}{J/K}$。（1）まわりは $\qty{-5.76}{J/K}$ で、合計0。（2）まわりは0で、合計 $+\qty{5.76}{J/K}$。
\phtashikame
（1）は可逆なので合計0、（2）は不可逆なので合計が正で、⑥と合います。\\
単位：mol×J/(mol·K)＝J/K。温度 $T$ は答えに入りません。\\
熱-15 では、同じ\qty{5.76}{J/K}を、分子の並び方の数から求めます。
\kotae 気体はどちらも\qty{5.76}{J/K}増える。合計は（1）0、（2）\qty{5.76}{J/K}
\end{phreidai}
\begin{jibun}
理想気体\qty{2}{mol}が、温度を変えずに自由膨張して、体積が3倍になりました。気体のエントロピーの変化を求めましょう。
\bkjans{$2\times8.314\times\ln3\approx\qty{18.3}{J/K}$ 増える}
\end{jibun}

\label{ph:熱-6:ex5}
\begin{phreidai}[やかんの水を温めるときのエントロピー]
水\qty{1.0}{kg}を、20\,℃から100\,℃まで温めます（わかさない）。水のエントロピーの変化を求めましょう。3けたで書くので、比熱は\qty{4.18}{kJ/(kg.K)}を使います。
\phbamen
やかんの水を温める場面です。温度が上がりながら熱を受けとります。
\phmushi
水の蒸発、やかんの熱容量を無視します。比熱は一定とします。
\phzu
$m=\qty{1.0}{kg}$、$c=\qty{4180}{J/(kg.K)}$、$T_1=\qty{293.15}{K}$、$T_2=\qty{373.15}{K}$。
\phriyu
温度が変わるので、$Q/T$ を1回では使えません。⑦の $mc\ln(T_2/T_1)$ を使います。
\phtateru
\[\Delta S=mc\ln\frac{T_2}{T_1}\]
\phtoku
\[\Delta S=1.0\times4180\times\ln\frac{373.15}{293.15}=4180\times0.2413\approx\qty{1.01e3}{J/K}\]
\phtashikame
はさみうち：受けとる熱 $Q=4180\times80=\qty{334400}{J}$ を、いちばん高い温度でわると\qty{896}{J/K}、いちばん低い温度でわると\qty{1141}{J/K}。答えはその間にあります。\\
間違えやすい所：℃のまま $\ln\tfrac{100}{20}$ とすると、約6700 J/K と大きくずれます（まちがえやすいところ）。
\kotae 約\qty{1.01e3}{J/K}増える
\end{phreidai}
\begin{jibun}
水\qty{2.0}{kg}を、10\,℃から60\,℃まで温めるときの、エントロピーの変化を3けたで求めましょう（比熱\qty{4.18}{kJ/(kg.K)}）。
\bkjans{$2.0\times4180\times\ln(333.15/283.15)$ より、約\qty{1.36e3}{J/K}増える}
\end{jibun}

\label{ph:熱-6:ex6}
\begin{phreidai}[カルノーサイクルを $T$-$S$ 線図で読む]
熱-5 例題8 のカルノーサイクル（理想気体\qty{1}{mol}、高温\qty{600}{K}、低温\qty{300}{K}、等温の膨張で体積が2倍）を $T$-$S$ 線図にかき、受けとる熱・捨てる熱・正味の仕事を、面積から求めましょう。
\phbamen
熱-5 で $p$-$V$ 図で調べたサイクルを、もう1つの図で見直します。
\phmushi
すべての変化が可逆とします（カルノーサイクル）。
\phzu
等温の膨張で $S$ が $\Delta S$ だけ増え、断熱で可逆な膨張では $S$ は変わらず、等温の圧縮で $\Delta S$ だけ減ります。図\ref{ph:熱-6:fig-ts}の（イ）の長方形です。
\phriyu
等温の膨張の $\Delta S$ は④の $nR\ln2$。面積は⑧のとおり、上の辺の下が $Q_{\mathrm H}$、下の辺の下が $Q_{\mathrm L}$、長方形が $W_{\text{出}}$ です。
\phtateru
\[Q_{\mathrm H}=T_{\mathrm H}\Delta S,\qquad Q_{\mathrm L}=T_{\mathrm L}\Delta S,\qquad W_{\text{出}}=(T_{\mathrm H}-T_{\mathrm L})\Delta S\]
\phtoku
$\Delta S=8.314\times\ln2\approx\qty{5.763}{J/K}$。$Q_{\mathrm H}=600\times5.763\approx\qty{3460}{J}$、$Q_{\mathrm L}=300\times5.763\approx\qty{1730}{J}$、$W_{\text{出}}\approx\qty{1730}{J}$。
\phtashikame
効率 $W_{\text{出}}/Q_{\mathrm H}=0.50=1-\tfrac{300}{600}$ で、カルノーの効率と同じ。\\
熱-5 例題8 で $p$-$V$ 図の4つの道の $\int p\,\dd V$ をたした仕事\qty{1730}{J}と同じです（⑧の「どちらの面積も正味の仕事」）。
\kotae $Q_{\mathrm H}$ 約\qty{3460}{J}、$Q_{\mathrm L}$ 約\qty{1730}{J}、正味の仕事 約\qty{1730}{J}
\end{phreidai}
\begin{jibun}
理想気体\qty{1}{mol}のカルノーサイクル（高温\qty{500}{K}、低温\qty{300}{K}、等温の膨張で体積が3倍）の、受けとる熱と正味の仕事を、$T$-$S$ 線図の面積から求めましょう。
\bkjans{$\Delta S=8.314\ln3\approx\qty{9.13}{J/K}$。$Q_{\mathrm H}\approx\qty{4570}{J}$、正味の仕事 $200\times9.13\approx\qty{1830}{J}$（効率0.40）}
\end{jibun}

\label{ph:熱-6:ex7}
\begin{phreidai}[熱のエクセルギー]
\qty{600}{K}の熱\qty{1000}{J}から、まわり（\qty{300}{K}）を使って、いちばん多く取り出せる仕事（エクセルギー）を求めましょう。
\phbamen
工場の高温の排ガスの熱を、発電に使えるか考える場面です。
\phmushi
熱の源の温度は一定（\qty{600}{K}）とし、機械はすべて可逆とします。
\phzu
$Q=\qty{1000}{J}$、$T=\qty{600}{K}$、$T_0=\qty{300}{K}$。
\phriyu
いちばん多く取り出せるのは、$T$ と $T_0$ の間の可逆な機関なので、⑨の $E_{\mathrm x}=Q(1-T_0/T)$ です。
\phtateru
\[E_{\mathrm x}=Q\Bigl(1-\frac{T_0}{T}\Bigr)\]
\phtoku
\[E_{\mathrm x}=1000\times\Bigl(1-\frac{300}{600}\Bigr)=\qty{500}{J}\]
\phtashikame
残りの\qty{500}{J}はまわりへ捨てる熱で、$\tfrac{500}{300}=\tfrac{1000}{600}$ とエントロピーがそろう（可逆）。\\
極端な場合：$T=T_0$ なら0。まわりと同じ温度の熱からは、仕事は取り出せません。
\kotae \qty{500}{J}
\end{phreidai}
\begin{jibun}
\qty{900}{K}の熱\qty{2000}{J}の、まわり\qty{300}{K}でのエクセルギーを求めましょう。
\bkjans{$2000\times(1-\tfrac{300}{900})\approx\qty{1.33e3}{J}$}
\end{jibun}

\label{ph:熱-6:ex8}
\begin{phreidai}[混ぜると失われる仕事]
例題1で、お湯と水を混ぜたとき、失われた仕事（エクセルギー）を求めましょう。まわりの温度は $T_0=\qty{293.15}{K}$（20\,℃）とします。
\phbamen
混ぜる前なら、お湯と水の温度の差で熱機関を動かして、仕事を取り出せたはずです。混ぜてしまうと、その機会が失われます。
\phmushi
例題1と同じです。混ぜる以外の損失は考えません。
\phzu
$S_{\mathrm{gen}}=\qty{36.4}{J/K}$（例題1、⑦の式で計算した値）、$T_0=\qty{293.15}{K}$。
\phriyu
混ぜたとき、まわりとは熱をやりとりしないので、生まれたエントロピーは例題1の $\Delta S$ です。⑨の「失われる仕事＝$T_0S_{\mathrm{gen}}$」を使います。
\phtateru
\[W_{\text{失}}=T_0\,S_{\mathrm{gen}}\]
\phtoku
\[W_{\text{失}}=293.15\times36.4\approx\qty{1.07e4}{J}\approx\qty{11}{kJ}\]
\phtashikame
大きさ：お湯がもっていた熱（$4200\times30=\qty{126}{kJ}$）の1割より少しだけ小さい。温度の差が小さい熱は、仕事にしにくい。\\
⑤の「分けるのにいる最小の仕事」も同じ\qty{11}{kJ}です。失ったものを取りもどすには、同じだけの仕事がいります。
\kotae 約\qty{11}{kJ}
\end{phreidai}
\begin{jibun}
じぶんでとこう1（60\,℃と20\,℃の水\qty{1.0}{kg}ずつ）で、混ぜたときに失われた仕事を求めましょう（$T_0=\qty{293.15}{K}$）。
\bkjans{$293.15\times17.2\approx\qty{5.0e3}{J}$（約\qty{5.0}{kJ}）}
\end{jibun}

\label{ph:熱-6:ex9}
\begin{phreidai}[部屋の熱は、どれだけの仕事に当たる？]
外の気温が0\,℃のとき、20\,℃の部屋に入れる熱のエクセルギーは、熱の何\,\%ですか。
\phbamen
冬の暖房です。電気ストーブは、電気（エクセルギー100\,\%）をそのまま熱にします。
\phmushi
部屋の温度・外の温度は一定とします。
\phzu
$T=\qty{293.15}{K}$（部屋）、$T_0=\qty{273.15}{K}$（外）。
\phriyu
部屋の熱がもつエクセルギーの割合は、⑨の $1-T_0/T$ です。
\phtateru
\[\frac{E_{\mathrm x}}{Q}=1-\frac{T_0}{T}\]
\phtoku
\[1-\frac{273.15}{293.15}=\frac{20}{293.15}\approx0.068\]
\phtashikame
部屋に入れる熱の\qty{6.8}{\%}分の仕事があれば、理想のヒートポンプで同じ暖房ができる（COP は最大で $1/0.068\approx15$）。\\
電気ストーブは、電気の93\,\%分のエクセルギーをむだにしています（図\ref{ph:熱-6:fig-exergy}）。だからヒートポンプのほうが得なのです（熱-5 ④）。
\kotae 約6.8\,\%
\end{phreidai}
\begin{jibun}
外の気温が $-5$\,℃、部屋が25\,℃のとき、部屋に入れる熱のエクセルギーの割合を求めましょう。
\bkjans{$1-(268.15/298.15)\approx0.10$（約10\,\%）}
\end{jibun}

\label{ph:熱-6:ex10}
\begin{phreidai}[断熱で可逆な変化から、ポアソンの式を導く]
理想気体が、断熱で可逆に体積を変えるとき、$pV^{\gamma}$ が一定になることを、エントロピーの式から導きましょう。
\phbamen
ゆっくり押し縮める、熱を通さないシリンダーの中の空気です（熱-3 ⑥と同じ場面）。
\phmushi
理想気体で、比熱は一定。摩擦や漏れはなく、熱の出入りもないとします。
\phzu
はじめ $(T_1,\ V_1,\ p_1)$、あと $(T_2,\ V_2,\ p_2)$。比熱比 $\gamma=c_p/c_v$。
\phriyu
断熱（$\delta Q=0$）で可逆なので $\dd S=\delta Q_{\mathrm{rev}}/T=0$、エントロピーは変わりません。⑦の理想気体の式で $\Delta s=0$ とおきます。
\phtateru
\[c_v\ln\frac{T_2}{T_1}+R_{\mathrm s}\ln\frac{V_2}{V_1}=0\]
\phtoku
$c_v$ でわって、$R_{\mathrm s}/c_v=\gamma-1$ を使うと、$\ln(T_2/T_1)+(\gamma-1)\ln(V_2/V_1)=0$、つまり $\ln\bigl(T_2V_2^{\,\gamma-1}\bigr)=\ln\bigl(T_1V_1^{\,\gamma-1}\bigr)$。だから $TV^{\gamma-1}$ が一定。$T=pV/(nR)$ を入れて、$pV^{\gamma}$ が一定。
\phtashikame
熱-3 ⑥の変数分離の答えと同じです。\\
$\gamma=1$（$R_{\mathrm s}=0$）とすると $T$ 一定になり、形の上では等温にもどります。
\kotae $pV^{\gamma}=$ 一定（$TV^{\gamma-1}=$ 一定）
\end{phreidai}
\begin{jibun}
例題10と同じ考えで、圧力で書いた式 $\Delta s=c_p\ln(T_2/T_1)-R_{\mathrm s}\ln(p_2/p_1)$ から、断熱で可逆なときに $T/p^{(\gamma-1)/\gamma}$ が一定になることを導きましょう。
\bkjans{$\Delta s=0$ から $\ln(T_2/T_1)=(R_{\mathrm s}/c_p)\ln(p_2/p_1)$。$R_{\mathrm s}/c_p=(c_p-c_v)/c_p=(\gamma-1)/\gamma$ なので、$T/p^{(\gamma-1)/\gamma}$ が一定}
\end{jibun}

\label{ph:熱-6:ex11}
\begin{phreidai}[摩擦の熱で生まれるエントロピー]
床の上で箱を滑らせて止めると、摩擦で\qty{100}{J}の熱が出て、\qty{300}{K}の床と空気に入りました。生まれたエントロピーを求めましょう。
\phbamen
図\ref{ph:熱-6:fig-rei}の（イ）の場面です。仕事（運動エネルギー）が熱になりました。
\phmushi
床と空気は大きく、温度は\qty{300}{K}のまま変わらないとします。
\phzu
$Q=\qty{100}{J}$、$T=\qty{300}{K}$。
\phriyu
はじめは熱の出入りがなく、まわり（\qty{300}{K}）がエントロピー $Q/T$ を受けとっただけです。これが生まれた分です。
\phtateru
\[S_{\mathrm{gen}}=\frac{Q}{T}\]
\phtoku
\[S_{\mathrm{gen}}=\frac{100}{300}\approx\qty{0.33}{J/K}\]
\phtashikame
失われた仕事は $T_0S_{\mathrm{gen}}=300\times0.333=\qty{100}{J}$ で、摩擦でなくなった運動エネルギーと同じです。\\
正なので、逆（床の熱で箱が動き出す）は起きません。
\kotae 約\qty{0.33}{J/K}
\end{phreidai}
\begin{jibun}
機械の軸受けで、1秒に\qty{50}{J}の摩擦の熱が出て、\qty{320}{K}の部品からまわりへ流れます。1分間に、部品で生まれるエントロピーを求めましょう。
\bkjans{$50\times60\div320\approx\qty{9.4}{J/K}$}
\end{jibun}

\label{ph:熱-6:ex12}
\begin{phreidai}[圧縮機を通る空気のエントロピー]
空気が、\qty{300}{K}・\qty{100}{kPa}から、\qty{600}{K}・\qty{500}{kPa}に圧縮されました。\qty{1}{kg}あたりのエントロピーの変化を求めましょう。$c_p=\qty{1005}{J/(kg.K)}$、$R_{\mathrm s}=\qty{287}{J/(kg.K)}$ とします。
\phbamen
圧縮機の入口と出口の空気です。熱が逃げないなら、可逆なら $\Delta s=0$ のはずです。
\phmushi
空気は理想気体で、比熱は一定。入口と出口の速さの差は無視します。
\phzu
$T_1=\qty{300}{K}$、$p_1=\qty{100}{kPa}$、$T_2=\qty{600}{K}$、$p_2=\qty{500}{kPa}$。
\phriyu
温度と圧力がわかっているので、⑦の圧力で書いた式を使います。
\phtateru
\[\Delta s=c_p\ln\frac{T_2}{T_1}-R_{\mathrm s}\ln\frac{p_2}{p_1}\]
\phtoku
\[\Delta s=1005\ln2-287\ln5\approx697-462\approx\qty{235}{J/(kg.K)}\]
\phtashikame
単位：J/(kg·K)。\\
正なので、断熱の圧縮機なら、中で摩擦などの不可逆な変化が起きている（可逆なら\qty{100}{kPa}→\qty{500}{kPa}で出口は約\qty{475}{K}のはず、熱-3 ⑥）。\\
圧力の比をとるので、kPa のままでよい（比には単位が残らない）。
\kotae 約\qty{235}{J/(kg.K)}増える
\end{phreidai}
\begin{jibun}
空気が\qty{290}{K}・\qty{100}{kPa}から\qty{580}{K}・\qty{800}{kPa}になりました。\qty{1}{kg}あたりのエントロピーの変化を求めましょう（$c_p=1005$、$R_{\mathrm s}=287$）。
\bkjans{$1005\ln2-287\ln8\approx697-597$ で、約\qty{100}{J/(kg.K)}増える}
\end{jibun}

\label{ph:熱-6:ex13}
\begin{phreidai}[本当の熱機関で、クラウジウスの不等式を確かめる]
ある熱機関は、1周で、\qty{600}{K}から熱\qty{1000}{J}を受けとり、\qty{300}{K}へ熱\qty{600}{J}を捨てます。$\oint\delta Q/T$ と、生まれるエントロピー、失われる仕事を求めましょう。
\phbamen
摩擦などのある、本当のエンジンです。カルノーの機関より効率が低いはずです。
\phmushi
熱の源の温度は一定。働く物は1周でもとにもどります。
\phzu
受けとる向きを正にして、$Q_1=\qty{1000}{J}$（$T_1=\qty{600}{K}$）、$Q_2=\qty{-600}{J}$（$T_2=\qty{300}{K}$）。取り出す仕事 $W_{\text{出}}=1000-600=\qty{400}{J}$。
\phriyu
⑥のクラウジウスの不等式で確かめます。働く物はもとにもどるので、生まれたエントロピーは、2つの熱の源のエントロピーの変化の合計です。
\phtateru
\[\oint\frac{\delta Q}{T}=\frac{Q_1}{T_1}+\frac{Q_2}{T_2},\qquad S_{\mathrm{gen}}=-\frac{Q_1}{T_1}-\frac{Q_2}{T_2}\]
\phtoku
$\oint\delta Q/T=\tfrac{1000}{600}-\tfrac{600}{300}=1.67-2.00\approx\qty{-0.33}{J/K}$。$S_{\mathrm{gen}}\approx\qty{0.33}{J/K}$、失われる仕事 $300\times0.333=\qty{100}{J}$。
\phtashikame
負なので、不等式を満たし、この機関はつくれる。効率0.40はカルノーの0.50より低い。\\
カルノーなら\qty{500}{J}の仕事が出るので、$500-400=\qty{100}{J}$ が失われる仕事と合います。\\
もし $\oint$ が正になる機関があったら、それは第二法則に反し、つくれません。
\kotae $\oint\delta Q/T\approx\qty{-0.33}{J/K}$、生まれるエントロピー約\qty{0.33}{J/K}、失われる仕事\qty{100}{J}
\end{phreidai}
\begin{jibun}
「\qty{500}{K}から\qty{1000}{J}を受けとり、\qty{450}{J}の仕事をして、\qty{300}{K}へ\qty{550}{J}を捨てる機関」はつくれますか。$\oint\delta Q/T$ で判断しましょう。
\bkjans{$\tfrac{1000}{500}-\tfrac{550}{300}\approx2.00-1.83=\qty{+0.17}{J/K}>0$ なので、つくれない（効率0.45がカルノーの0.40をこえている）}
\end{jibun}

%% ======================== まちがえやすいところ ========================
\Needspace{11\baselineskip}\bksection{tsumazuki}{まちがえやすいところをチェック}
\begin{tsumazuki}{温度は℃のまま？}
\bkng{20\,℃から100\,℃に温める水の $\Delta S$ は、$mc\ln\tfrac{100}{20}$。}
\bkok{絶対温度で、$mc\ln(373.15/293.15)$。}
$\Delta S=Q/T$ も $\ln(T_2/T_1)$ も、温度の比をとる式です。℃で比をとると、0\,℃でわり算ができなくなり、答えも6倍以上ずれます（例題5）。
\end{tsumazuki}

\begin{tsumazuki}{熱が入らなければ、エントロピーは変わらない？}
\bkng{自由膨張は $Q=0$ なので、$\Delta S=0$。}
\bkok{自由膨張でも $\Delta S=nR\ln(V_2/V_1)>0$。}
$\Delta S=\int\delta Q/T$ が使えるのは、可逆な道だけです。不可逆な変化では、同じはじめとおわりを結ぶ可逆な道を考えて計算します（④）。「断熱なら $\Delta S=0$」は「断熱で\textbf{可逆}なら」です。
\end{tsumazuki}

\begin{tsumazuki}{エントロピーが減ったら、第二法則に反する？}
\bkng{水が凍ると、水のエントロピーが減るので、第二法則に反する。}
\bkok{減ってよい。まわりのエントロピーが、それ以上に増えている。}
増大の法則は、\textbf{孤立した系}（まわりまで全部入れた合計）についての法則です。凍る水は、熱をまわりに捨て、まわりのほうが温度が低いので、まわりの増えが大きくなります（章末問題2）。
\end{tsumazuki}

\begin{tsumazuki}{$T$-$S$ 線図の面積は、いつも熱？}
\bkng{不可逆な変化でも、$T$-$S$ 線図の線の下の面積は、受けとった熱。}
\bkok{面積が熱になるのは、可逆な変化だけ。}
不可逆な変化では、熱が入らなくても $S$ が増えます（自由膨張）。$T\,\dd S=\delta Q$ は、$\dd S=\delta Q_{\mathrm{rev}}/T$ から来た式です（⑧）。
\end{tsumazuki}

\begin{tsumazuki}{エネルギーが残っていれば、また使える？}
\bkng{ぬるま湯にも、混ぜる前と同じエネルギーがあるので、同じだけ仕事に使える。}
\bkok{エネルギーは同じでも、エクセルギー（仕事にできる分）は減っている。}
エネルギーは保存しますが、エクセルギーは不可逆な変化で失われます（⑨、例題8）。「エネルギーを節約する」とは、正確には「エクセルギーをむだにしない」ことです。
\end{tsumazuki}

\begin{tsumazuki}{$T$-$S$ 線図の面積と $p$-$V$ 図の面積}
\bkng{$T$-$S$ 線図の、閉じた道の下の面積が、正味の仕事。}
\bkok{正味の仕事は、閉じた道が\textbf{囲む}面積。道の下の面積は、その区間で受けとる（捨てる）熱。}
$T$-$S$ 線図でも $p$-$V$ 図でも、上の道の下の面積から下の道の下の面積を引いた残り（囲む面積）が、1周の正味の仕事です（⑧、例題6）。
\end{tsumazuki}

%% ======================== 工学ではこう使う ========================
\Needspace{18\baselineskip}\bksection{phgear}{工学ではこう使う}
\begin{phkougaku}[損失がどこで生まれるかを数える]
発電所・エンジン・冷凍機の設計では、機械の部分ごとに、生まれるエントロピー $S_{\mathrm{gen}}$ を計算します。それに $T_0$ をかけた $T_0S_{\mathrm{gen}}$ が、その部分で失われた仕事です。エネルギーのバランス（第一法則）だけでは、「どこを直せば効率が上がるか」はわかりません。たとえば、火力発電所では、エネルギーで見ると「捨てる熱がいちばん多いのは復水器（熱-0 ③）」ですが、エクセルギーで見ると、復水器で捨てる熱は温度が低いので、失われる仕事は小さく、いちばん大きく失われているのは、燃料が燃えて高い温度の熱になる所（ボイラー）だとわかります。

\phboxhead{温度の差を小さくする}
熱交換器では、熱い流れと冷たい流れの温度の差が大きいほど、早く熱が伝わり、装置は小さくてすみます。でも、例題3のとおり、温度の差が大きいほど、多くのエントロピーが生まれます。装置の大きさ（費用）と、失われる仕事（運転の費用）のつり合いをとるのが、設計です（熱-13）。

\begin{phketa}[エントロピーの大きさ]
\phketaline{水\qty{1}{kg}を20\,℃から100\,℃に}{約\qty{1.0}{kJ/K}}
\phketaline{氷\qty{1}{kg}がとける（0\,℃）}{約\qty{1.2}{kJ/K}}
\phketaline{水\qty{1}{kg}がわく（100\,℃）}{約\qty{6.0}{kJ/K}}
\phketaline{理想気体\qty{1}{mol}が体積2倍}{\qty{5.76}{J/K}}
\phboxhead{答えが桁で合っているか}
\qty{1}{kg}の物の温度・状態の変化なら、エントロピーの変化は数百〜数千 J/K です。$\Delta S$ が数十万 J/K になったら、℃のまま比をとったか、kJ と J を間違えています。失われる仕事 $T_0S_{\mathrm{gen}}$ は、$T_0$ が約\qty{300}{K}なので、$S_{\mathrm{gen}}$〔J/K〕の約300倍〔J〕です。
\end{phketa}
\end{phkougaku}
\phstation{工業熱力学,エネルギー工学,伝熱工学,化学工学}{エントロピー、エントロピー生成、エクセルギー（有効エネルギー）、有効エネルギー損失、等エントロピー変化、$T$-$s$ 線図}

\begin{column}{「エントロピー」という名前}
エントロピーということばは、ドイツの物理学者クラウジウス（1822〜1888年）が、1865年の論文でつけました。ギリシャ語の「変化（トロペー）」をもとに、「エネルギー」と似た形になるように選んだ、と本人が書いています。クラウジウスは、論文の最後で、熱力学の2つの法則を「宇宙のエネルギーは一定である。宇宙のエントロピーは最大に向かう」とまとめました。この章の⑥の不等式も、クラウジウスが導いたものです。
\end{column}

\begin{column}{エクセルギーと、失われる仕事}
「失われる仕事は、まわりの温度×生まれるエントロピー」という関係は、フランスのグイが1889年に、スイスで蒸気タービンを研究したストドラが1900年代のはじめに、それぞれ示しました。「エクセルギー」という名前は、1956年にスロベニアのラントが、「仕事として取り出せる能力」を表すことばとしてつくりました。日本の工学の本では「有効エネルギー」ともいいます。
\end{column}

\begin{column}{マクスウェルの悪魔}
1867年ごろ、マクスウェルは、こんな思考実験を考えました。箱を仕切りで2つに分け、仕切りに小さな戸をつけます。戸の番をする小さな「悪魔」が、速い分子だけを右へ、遅い分子だけを左へ通すと、仕事をしないで、右が熱く、左が冷たくなり、エントロピーが減るように見えます。この「悪魔」のなぞは、長い間の議論になりました。今では、悪魔が分子の速さを測って覚えておき、その記録を消すときに、まわりのエントロピーが増えるので、全体では第二法則は破れない、と考えられています（ランダウアーとベネットの考え。記録を1ビット消すのにいる最小のエネルギーは熱-15）。情報とエントロピーのこの関係は、2010年代に、とても小さな粒を使った実験でも確かめられています。
\end{column}

%% ======================== 章末問題 ========================
\bkendproblems
\bkproblemgroup{基本}{3}
\bkproblem 熱\qty{600}{J}を、（1）\qty{300}{K}の大きな水そうが受けとるとき、（2）\qty{600}{K}の大きな水そうが受けとるときの、それぞれのエントロピーの変化を求めましょう。
\bkpans{（1）\nobreak \qty{2.0}{J/K}。（2）\nobreak \qty{1.0}{J/K}}
\bkproblem 0\,℃の水\qty{0.20}{kg}が、$-10$\,℃の冷凍庫の中で凍って0\,℃の氷になりました（融解熱\qty{334}{kJ/kg}）。水のエントロピーの変化、冷凍庫の中の空気（大きく、温度は $-10$\,℃のまま）のエントロピーの変化、その合計を求めましょう。
\bkpans{水 約 $\qty{-245}{J/K}$、冷凍庫 約 $66800\div263.15\approx\qty{+254}{J/K}$、合計 約 $\qty{+9}{J/K}$（増えている）}
\bkproblem 水\qty{0.50}{kg}を、20\,℃から80\,℃に温めるときのエントロピーの変化を求めましょう（$c=\qty{4.2}{kJ/(kg.K)}$）。
\bkpans{$0.50\times4200\times\ln(353.15/293.15)$ より、約\qty{3.9e2}{J/K}}
\bkproblem 同じ温度・同じ圧力の、窒素\qty{1}{mol}と酸素\qty{1}{mol}を、仕切りの左右に入れておき、仕切りをとって、混ぜました。全体のエントロピーの変化を求めましょう。また、理想気体\qty{2}{mol}が自由膨張して体積が2倍になるときとくらべましょう。
\bkpans{$2R\ln2\approx\qty{11.5}{J/K}$。\qty{2}{mol}の自由膨張（$2R\ln2$）と同じ}
\bkproblem まわりの温度を\qty{300}{K}として、（1）\qty{800}{K}の熱\qty{1000}{J}、（2）\qty{400}{K}の熱\qty{1000}{J}のエクセルギーを求めましょう。
\bkpans{（1）\nobreak \qty{625}{J}。（2）\nobreak \qty{250}{J}}
\bkproblem ある熱機関は、1周で\qty{500}{K}から熱\qty{1000}{J}を受けとり、\qty{300}{J}の仕事をして、\qty{300}{K}へ熱を捨てます。$\oint\delta Q/T$、生まれるエントロピー、失われる仕事（$T_0=\qty{300}{K}$）を求めましょう。
\bkpans{捨てる熱\qty{700}{J}。$\oint\delta Q/T=2.00-2.33\approx\qty{-0.33}{J/K}$。生まれるエントロピー約\qty{0.33}{J/K}、失われる仕事\qty{100}{J}（カルノーなら\qty{400}{J}）}
\bkproblem 空気\qty{1}{kg}を、\qty{300}{K}から\qty{450}{K}に、（1）圧力一定で、（2）体積一定で温めるときのエントロピーの変化を求めましょう（$c_p=\qty{1005}{J/(kg.K)}$、$c_v=\qty{718}{J/(kg.K)}$）。
\bkpans{（1）\nobreak $1005\ln1.5\approx\qty{408}{J/(kg.K)}$。（2）\nobreak $718\ln1.5\approx\qty{291}{J/(kg.K)}$}
\bkproblem 理想気体のカルノーサイクルを $T$-$S$ 線図にかくと、高温\qty{800}{K}、低温\qty{300}{K}、横の辺の長さ（エントロピーの変化）が\qty{3.0}{J/K}の長方形になりました。受けとる熱、捨てる熱、正味の仕事、効率を、面積から求めましょう。
\bkpans{受けとる熱 $800\times3.0=\qty{2400}{J}$、捨てる熱 $300\times3.0=\qty{900}{J}$、正味の仕事\qty{1500}{J}、効率 $1500\div2400=0.625$（$1-300/800$ と同じ）}
\bkproblemgroup{チャレンジ}{5}
\bkproblem 銅のかたまり\qty{1.0}{kg}が2つあり、温度は\qty{400}{K}と\qty{300}{K}です（比熱\qty{385}{J/(kg.K)}）。（1）2つをくっつけて温度をそろえたときの、全体のエントロピーの変化を求めましょう。（2）くっつける代わりに、2つの間で可逆な熱機関を動かして温度をそろえると、全体のエントロピーは変わりません。このことから、最後の温度と、取り出せる仕事を求めましょう。
\bkpans{（1）\nobreak \qty{350}{K}になり、$385\bigl(\ln\tfrac{350}{400}+\ln\tfrac{350}{300}\bigr)\approx\qty{7.9}{J/K}$。（2）\nobreak $\ln(T_{\mathrm f}/400)+\ln(T_{\mathrm f}/300)=0$ から $T_{\mathrm f}=\sqrt{400\times300}\approx\qty{346}{K}$。仕事は $385\times(400+300-2\times346.4)$ より、約\qty{2.8}{kJ}}
\bkproblem $T$-$S$ 線図で、次の3本の線でできる三角形のサイクルがあります。点 $(S_1,\ \qty{500}{K})$ から $(S_1+\qty{2}{J/K},\ \qty{500}{K})$ へ等温で、そこから点 $(S_1,\ \qty{300}{K})$ へまっすぐ、そこから $(S_1,\ \qty{500}{K})$ へ断熱で可逆にもどります。受けとる熱、捨てる熱、正味の仕事、効率を求め、同じ温度の間のカルノーの効率とくらべましょう。
\par\smallskip\noindent\begin{center}
\begin{tikzpicture}[x=1mm, y=1mm]
  \draw[bk arrow] (0,0) -- (52,0) node[right, font=\sffamily\footnotesize] {$S$};
  \draw[bk arrow] (0,0) -- (0,34) node[above, font=\sffamily\footnotesize] {$T$〔K〕};
  \fill[bkOrange!20] (12,28) -- (36,28) -- (12,16) -- cycle;
  \draw[bkInk, line width=1.2pt, -{Stealth[length=2mm]}] (12,28) -- (36,28);
  \draw[bkInk, line width=1.2pt, -{Stealth[length=2mm]}] (36,28) -- (12.5,16.25);
  \draw[bkInk, line width=1.2pt, -{Stealth[length=2mm]}] (12,16) -- (12,27.5);
  \draw[bkGray, dashed, line width=0.4pt] (0,28) -- (12,28) (0,16) -- (12,16) (36,0) -- (36,28) (12,0) -- (12,16);
  \node[font=\sffamily\footnotesize, anchor=east] at (-1,28) {500};
  \node[font=\sffamily\footnotesize, anchor=east] at (-1,16) {300};
  \node[font=\sffamily\footnotesize, anchor=north] at (12,-1) {$S_1$};
  \node[font=\sffamily\footnotesize, anchor=north] at (36,-1) {$S_1+2$};
\end{tikzpicture}
\end{center}
\bkpans{受けとる熱 $500\times2=\qty{1000}{J}$、捨てる熱（ななめの線の下の台形）$(500+300)\div2\times2=\qty{800}{J}$、正味の仕事（三角形）\qty{200}{J}、効率0.20（カルノーは0.40）}
\bkproblem 外の気温が0\,℃のとき、20\,℃の部屋を\qty{2.0}{kW}の熱で暖めます。（1）理想のヒートポンプで暖めるときの、いちばん小さい電力、（2）電気ストーブ（電力\qty{2.0}{kW}）で暖めるとき、失われる仕事率を求めましょう。
\bkpans{（1）\nobreak $2000\times0.068\approx\qty{1.4e2}{W}$（COP 約15）。（2）\nobreak 約 $2000-136\approx\qty{1.9}{kW}$}
\bkproblem 0\,℃の氷\qty{1.0}{kg}を、20\,℃の大きな部屋に置いておくと、とけて、やがて20\,℃の水になりました（融解熱\qty{334}{kJ/kg}、水の比熱\qty{4.2}{kJ/(kg.K)}）。（1）氷（水）のエントロピーの変化、（2）部屋の空気のエントロピーの変化、（3）生まれたエントロピーと、失われた仕事（$T_0=\qty{293.15}{K}$）を求めましょう。
\bkpans{（1）\nobreak とける分 $334000\div273.15\approx1223$、温まる分 $4200\ln(293.15/273.15)\approx297$ で、約\qty{1.52e3}{J/K}。（2）\nobreak 部屋が出した熱は $334000+84000=\qty{418000}{J}$ なので、$-418000\div293.15$ より、約 $\qty{-1.43e3}{J/K}$。（3）\nobreak 約\qty{94}{J/K}、失われた仕事 約\qty{27}{kJ}}

\begin{furikaeri}
\begin{checklist}
\item 熱÷温度 $\Delta S=Q/T$ を数えて、もどらない向きを言える。\bkfurref{ph:熱-6:sec2}
\item 水・氷・気体の $\Delta S$ を計算し、$T$-$s$ 線図の面積を読める。\bkfurref{ph:熱-6:sec7}
\item 生まれるエントロピーとエクセルギーで、機械の損失を数えられる。\bkfurref{ph:熱-6:sec9}
\end{checklist}
\end{furikaeri}
\bkfukidashi{mina}{エネルギーはなくならないのに、使える分は減っていくんだね。お湯と水を混ぜるだけで、\qty{11}{kJ}も仕事のチャンスを捨てていたなんて。}
